\documentclass[review,3p]{elsarticle}

%% ============================================================
%% PAPER 1: TOPOLOGY ABSTRACTION & PHYSICS FIDELITY EFFECTS
%% Target Journal: Applied Energy
%% Version: 7.3 — Cleaned formatting
%% ============================================================

\usepackage[utf8]{inputenc}
\usepackage[T1]{fontenc}
\usepackage{amsmath,amssymb,amsfonts}
\usepackage{graphicx}
\usepackage{booktabs}
\usepackage{multirow}
\usepackage{array}
\usepackage{tabularx}
\usepackage{siunitx}
\usepackage{xcolor}
\usepackage{hyperref}
\usepackage{cleveref}
\usepackage{lineno}
\usepackage{float}
\usepackage{subfig}
\usepackage{enumitem}
\usepackage[version=4]{mhchem}

\modulolinenumbers[5]

\sisetup{
  per-mode=symbol,
  range-phrase=--,
  range-units=single
}

\hypersetup{
  colorlinks=true,
  linkcolor=blue,
  citecolor=blue,
  urlcolor=blue
}

%% --- Custom commands ---
\newcommand{\placeholder}[1]{\textcolor{red}{[#1]}}
\newcommand{\Lthree}{\mathrm{L3}}
\newcommand{\Lplus}{\mathrm{L3}^{+}}
\newcommand{\Lnl}{\mathrm{L3}^{\mathrm{NL}}}
\newcommand{\result}[1]{\textcolor{blue}{[\textsc{#1}]}}
\newcommand{\figref}[1]{\textcolor{teal}{[FIG:#1]}}
\newcommand{\todo}[1]{\textcolor{orange}{\textbf{TODO:} #1}}

%% --- Column type for centered, fixed-width ---
\newcolumntype{C}[1]{>{\centering\arraybackslash}m{#1}}
\newcolumntype{L}[1]{>{\raggedright\arraybackslash}m{#1}}

\journal{Applied Energy}

\begin{document}

\begin{frontmatter}

\title{Topology Abstraction and Physics Fidelity Effects on Dispatch
Optimization of Electrified District Heating Networks:\\
From Copperplate MILP to Quadratic Network Models}

\author[IPA]{Lukas Ruess\corref{cor1}}
\ead{lukas.ruess@ipa.fraunhofer.de}
\ead{lukas.ruess@eep.uni-stuttgart.de}
\author[EEP]{Alexander Sauer}
\author[e-con]{Stefan Roth}


\address[IPA]{Fraunhofer-Institut für Produktionstechnik und Automatisierung, Stuttgart, Germany}
\address[EEP]{Institut für Energieeffizienz in der Produktion, Stuttgart, Germany}
\address[e-con]{e-con AG, Memmingen, Germany}


\begin{highlights}
  \item First controlled comparison isolating topology, physics, and
        linearization effects in DH dispatch
  \item Physics inclusion (losses) dominates spatial refinement by
        factor~4 in cost impact
  \item Piecewise-linear MILP deviates $<$1\,\% from
        quadratic reference---adequate for planning
  \item Physics inclusion gap (13\,\%) exceeds spatial refinement gap
        (2.5\,\%) and linearization error (0.35\,\%) by order of magnitude
  \item Decision criteria link pipe length and demand heterogeneity
        to recommended model fidelity level
\end{highlights}

%% ============================================================
%% ABSTRACT
%% ============================================================
\begin{abstract}
District heating dispatch optimization requires choosing both a
topology representation level and a physics fidelity level, yet no
controlled quantitative guidance exists for either decision. Prior
comparisons confound these dimensions simultaneously.

We present the first study that isolates topology abstraction,
physics fidelity, and linearization error as independent experimental
variables across five model variants: L1~(copperplate), L2~(7-zone
MILP), L3~(15-node MILP, basic physics), $\Lplus$~(15-node MILP,
extended physics), and $\Lnl$~(15-node MIQCP reference). The
framework is validated against a German industrial heat network and
assessed across 36~synthetic configurations.

Results establish a clear effect hierarchy: the L1$\to$L2 transition
(physics inclusion) causes 12\,\% cost underestimation and
63\,t\,\ce{CO2}/year; the full L1$\to$L3 gap reaches
74\,t\,\ce{CO2}/year; L3$\to\Lplus$ (extended physics) adds
$0.11$\,\% through pumping and transport delay; and the
$\Lplus\to\Lnl$ linearization error is ${\leq}\,0.35$\,\%,
validating MILP for planning purposes. These findings yield
actionable selection criteria linking pipe length, demand
heterogeneity, and transport delay to the minimum warranted model
level.
\end{abstract}

\begin{keyword}
  District heating \sep
  Mixed-integer linear programming \sep
  Mixed-integer quadratically constrained programming \sep
  Topology abstraction \sep
  Thermal storage dispatch \sep
  Heat pump optimization \sep
  Network modeling fidelity \sep
  Linearization error
\end{keyword}

\begin{dataavailability}
The optimization model source code (Python/Pyomo), all YAML
configuration files defining the five model variants, and the
post-processing scripts for reproducing all figures and tables are
publicly available on Zenodo at
\url{https://doi.org/10.5281/zenodo.20394195} under an MIT license.
The input time series (electricity prices, heat demand profiles,
waste-heat source temperatures) and pre-upgrade monitoring data are
subject to a non-disclosure agreement with the network operator and
cannot be shared publicly. Anonymized summary statistics sufficient
to verify key results are included in the repository. Researchers
seeking access to the full input dataset may contact the
corresponding author.
\end{dataavailability}

\begin{credit}
\textbf{Lukas Ruess:} Conceptualization, Methodology, Software,
Validation, Formal analysis, Investigation, Data curation,
Writing~-- original draft, Visualization.
\textbf{Alexander Sauer:} Supervision, Writing~-- review \& editing,
Funding acquisition.
\textbf{Stefan Roth:} Resources, Data curation.
\end{credit}

\end{frontmatter}

\linenumbers

%% ============================================================
%% 1. INTRODUCTION
%% ============================================================
\section{Introduction}
\label{sec:introduction}

\subsection{Motivation and problem statement}
\label{sec:motivation}

Space and process heat account for approximately 50\,\% of final
energy consumption in the European Union, of which roughly 60\,\% is
still supplied by fossil fuels~\cite{EUheating2016}. District
heating~(DH) networks offer a systemic decarbonization pathway
through sector coupling: large-scale heat pumps, combined heat and
power~(CHP) units, thermal energy storage, and electrode boilers can
be coordinated within a single network~\cite{Lund2014_4GDH,Werner2017}.
Planning and operational optimization of such multi-source networks
increasingly relies on mathematical programming due to its optimality
guarantees and tractable handling of commitment decisions.

Every optimization-based planning study confronts two fundamental and
interrelated modeling choices:
%
\begin{enumerate}[nosep]
  \item \textbf{Topology abstraction:} the level of spatial detail at
        which the network is represented---from a single-bus
        copperplate to a pipe-by-pipe graph.
  \item \textbf{Physics fidelity:} the set of physical phenomena
        modeled---from zero-loss energy balances to
        pressure-drop-coupled, temperature-propagating, delay-aware
        thermo-hydraulic models.
\end{enumerate}

At one extreme, \emph{copperplate} models aggregate all components
and demands to a single bus, ignoring spatial structure and thermal
losses entirely. Wirtz et~al.~\cite{Wirtz2020} explicitly acknowledge
this limitation: ``the validity of [the global energy balance]
assumption depends on the network topology, size and the operation
strategy,'' yet no quantification of the resulting error has been
published. At the other extreme, \emph{full thermo-hydraulic} models
resolve the physical pipe-by-pipe network, capturing thermal losses
(5--15\,\% of annual heat
supply~\cite{Frederiksen2013,Talebi2016}), pressure-dependent
pumping costs, temperature propagation, and transport delays arising
from finite fluid velocity. The practical consequences of these
choices affect cost estimation, dispatch quality, and \ce{CO2}
accounting simultaneously.

A further complication arises from the \emph{mathematical
formulation}: extended thermo-hydraulic physics introduce nonlinear
and bilinear terms (pressure drop is quadratic in flow;
temperature losses involve flow--temperature products) that are
incompatible with standard MILP solvers. Practitioners must choose
between linearization (preserving MILP tractability but introducing
approximation error) and nonlinear/quadratic formulations
(preserving physical accuracy but potentially sacrificing solve
time). The magnitude of the linearization error in a DH dispatch
context has not been quantified.

\subsection{Research gap}
\label{sec:gap}

The core problem is \emph{multi-dimensional confounding}: published
comparisons change multiple modeling dimensions simultaneously. For
instance, comparing the full nonlinear pipe model with
temperature-dependent COP used by Vesterlund et
al.~\cite{Vesterlund2017} against the copperplate formulation of
Gabrielli et al.~\cite{Gabrielli2018} conflates topology effects
with (i)~NLP vs.\ constant-loss formulations,
(ii)~temperature-dependent vs.\ constant COP,
(iii)~the presence or absence of hydraulic coupling, and
(iv)~differing temporal resolutions. The same pattern appears
systematically across the literature (\Cref{tab:prior_work}).

Four specific gaps persist:
%
\begin{enumerate}[nosep]
  \item No study isolates topology abstraction as the sole
        experimental variable while holding physics constant.
  \item No study isolates the physics fidelity effect (adding
        pressure drop, temperature propagation, transport delay)
        while holding topology constant.
  \item The linearization error introduced by MILP-compatible
        approximations of thermo-hydraulic physics has not been
        quantified against a quadratic reference in a DH dispatch
        context.
  \item No quantitative guidance exists for joint topology and
        physics selection as a function of observable network
        properties.
\end{enumerate}

\subsection{Contributions and research questions}
\label{sec:contributions}

This paper addresses these gaps through five contributions:
%
\begin{enumerate}
  \item \textbf{Five-level taxonomy.} A reusable classification
        framework for DH modeling fidelity---L1 (copperplate), L2
        (simplified multi-node), L3 (detailed MILP, basic physics),
        $\Lplus$ (detailed MILP, extended physics), $\Lnl$ (detailed
        MIQCP, extended physics without linearization)---enabling
        controlled comparison across both topology and physics
        dimensions (\Cref{sec:topology_def}).

  \item \textbf{First controlled topology comparison.} Three topology
        levels (L1/L2/L3) within a unified MILP framework where L2
        and L3 share identical physics---isolating topology as the
        sole experimental variable---validated against measured
        industrial operational data (\Cref{sec:results_cost}).

  \item \textbf{First controlled physics fidelity comparison.}
        L3 vs.\ $\Lplus$ isolates the effect of adding pressure drop,
        temperature-dependent loss propagation, and transport delay
        while holding topology constant
        (\Cref{sec:results_extended}).

  \item \textbf{Linearization error quantification.}
        $\Lplus$ (MILP) vs.\ $\Lnl$ (MIQCP) quantifies the cost of
        linearization (\Cref{sec:results_linearization}).

  \item \textbf{Actionable decision criteria.} Parametric analysis
        across 36~synthetic configurations yielding joint selection
        rules based on observable network properties
        (\Cref{sec:results_generalizability}).
\end{enumerate}

These contributions answer four research questions:
%
\begin{description}[style=nextline, nosep, leftmargin=1.5em]
  \item[RQ1 --- Topology effect:]
    How much does topology abstraction change annual operational cost
    and \ce{CO2} emissions when physics are held constant?
    (L1 vs.\ L2 vs.\ L3)

  \item[RQ2 --- Physics fidelity effect:]
    How much do pressure drop, temperature-dependent losses, and
    transport delay change dispatch and cost when topology is held
    constant? (L3 vs.\ $\Lplus$)

  \item[RQ3 --- Linearization error:]
    How large is the cost deviation introduced by MILP-compatible
    linearization compared to a quadratic reference?
    ($\Lplus$ vs.\ $\Lnl$)

  \item[RQ4 --- Generalizability:]
    Which network properties determine whether detailed topology
    and/or extended physics representation is warranted?
\end{description}

%% ============================================================
%% 2. RELATED WORK
%% ============================================================
\section{Related work}
\label{sec:related_work}

\subsection{MILP optimization of district heating systems}
\label{sec:rw_milp}

MILP has become the dominant methodology for operational optimization
of DH systems~\cite{Fang2016,Guelpa2019review}. Several open-source
frameworks support multi-carrier energy system optimization:
oemof~\cite{Hilpert2018_oemof}, PyPSA~\cite{Brown2018_pypsa},
FINE~\cite{Welder2018_FINE}, Calliope~\cite{Pfenninger2018_calliope},
and urbs~\cite{Dorfner2016_urbs}. Despite their diversity, these
frameworks share common limitations for detailed DH analysis: thermal
losses are typically time-invariant, heat pump COP is constant or
seasonal, hydraulic coupling is absent, and network topology is either
fully aggregated or fixed at a single detail level.

In the multi-energy system~(MES) literature,
Mancarella~\cite{Mancarella2014} established the theoretical basis for
multi-energy hub modeling, while Gabrielli et
al.~\cite{Gabrielli2018} formulated a MILP for distributed MES
including thermal and electrical storage with hourly resolution. These
studies optimize component dispatch and sizing but treat the thermal
network as a single aggregated efficiency factor---a simplification
whose cost this paper quantifies.

\subsection{Topology representation and spatial scale}
\label{sec:rw_topology}

The DH modeling literature exhibits a continuous spectrum of spatial
aggregation.

\paragraph{Single-node aggregation (L1-equivalent)}
Standard in national-scale energy system models where heating is one
of many sectors~\cite{Connolly2014_EnergyPLAN,Brown2018_pypsa,
Lund2017_SmartEnergy}. No spatial structure or thermal losses. Wirtz
et~al.~\cite{Wirtz2020} apply this approach to 5th-generation
bidirectional low-temperature networks, noting that residual building
demands are assumed to be covered ``perfectly'' by the energy
hub---an assumption this paper quantifies.

\paragraph{Zone-based models (L2-equivalent)}
Demand aggregated into geographic zones (3--15 nodes) connected by
simplified pipe
segments~\cite{Wirtz2020,Mertz2016,Buffa2019}. Bünning et
al.~\cite{Buenning2020} demonstrate that zone-level aggregation
captures 85--95\,\% of dynamic temperature variation.

\paragraph{Full-graph models (L3-equivalent)}
Node-by-node representation with pipe-by-pipe parameters from GIS
data~\cite{Vesterlund2017,Schweiger2017,vanderHeijde2017,
Blommaert2020}. Giraud et~al.~\cite{Giraud2015} show tractability up
to ${\approx}\,100$~nodes for hourly annual horizons.

Recent work on spatial-scale effects in DH
planning~\cite{Guo2025} demonstrates that the geographic scope of a
network (campus, district, or city-wide) fundamentally determines the
optimal system configuration. Our demand heterogeneity index
(HI, \Cref{sec:model_hi}) provides a complementary scalar metric:
while clustering validity indices determine \emph{how} to group
consumers, HI quantifies \emph{whether} grouping introduces
significant error by measuring load concentration independent of
spatial assignment.

\subsection{Thermo-hydraulic network models}
\label{sec:rw_thermohydraulic}

Beyond topology, the physics fidelity of pipe models varies
substantially.

\paragraph{Steady-state thermal losses}
The simplest physically meaningful model, where pipe losses follow
from the fluid--ground temperature difference and pipe insulation
properties~\cite{Paalsson1999,Frederiksen2013}. MILP-compatible when
temperatures are fixed parameters.

\paragraph{Pressure drop and pumping}
The Darcy--Weisbach equation relates pressure drop to the square of
flow velocity:
$\Delta p = f_D \frac{L}{D} \frac{\rho v^2}{2}$, introducing a
quadratic nonlinearity. Pumping power is proportional to
$\dot{m} \cdot \Delta p$, creating a cubic term in mass flow. Several
authors have addressed this through piecewise
linearization~\cite{Bordin2016,Mertz2016} or second-order cone
relaxations~\cite{Blommaert2020}. Hari et~al.~\cite{Hari2024} provide
a comprehensive thermal network flow optimization~(TNFO) framework
that integrates Darcy--Weisbach pressure drop with explicit pump
placement at multiple network locations. For hot-water networks
operating below \SI{100}{\celsius}, the incompressible-fluid
assumption decouples pressure drop from
temperature~\cite{Hari2024}. Reported pumping costs range from
1--5\,\% of thermal energy costs depending on network length and flow
velocity~\cite{Frederiksen2013}.

\paragraph{Temperature propagation}
In physical networks, heat losses along a pipe reduce the temperature
at downstream nodes. The temperature drop along a pipe of length~$L$
under steady-state conditions follows~\cite{vanderHeijde2017}:
%
\begin{equation*}
  T_{\mathrm{out}} = T_{\mathrm{ground}}
    + (T_{\mathrm{in}} - T_{\mathrm{ground}})
      \exp\!\left(\frac{-UL}{\dot{m}\,c_p}\right),
\end{equation*}
%
introducing a nonlinear (exponential) dependence on mass flow.
Maldonado et~al.~\cite{Maldonado2024} validate this model against
measured data from an industrial network (Lampoldshausen), achieving
supply temperature MAE below \SI{0.5}{\celsius} after
branch-sequential calibration---a methodology adopted in this work.

\paragraph{Transport delay}
Heat injected at the source reaches downstream consumers after a
delay $\tau = L/v$. For typical DH networks, transport times range
from minutes (short branches) to 1--3\,hours (long trunk
lines)~\cite{vanderHeijde2017,Giraud2015}. The fixed-delay
approximation is most accurate when flow variations are moderate
($\pm 30\,\%$ of nominal); Hari et~al.~\cite{Hari2024} show that
under contingency conditions (200\,\% load increase), flow rates can
vary by ${>}\,40\,\%$, potentially invalidating fixed-delay
assumptions.

\paragraph{Nonlinear and quadratic formulations}
Vesterlund et~al.~\cite{Vesterlund2017} formulated a full NLP for
multi-source DH optimization. Blommaert et~al.~\cite{Blommaert2020}
used adjoint-based NLP for topology design. Mixed-integer
quadratically constrained programs~(MIQCP) offer a middle ground:
quadratic constraints are handled natively by modern solvers (Gurobi,
CPLEX) without piecewise linearization, while integer variables retain
unit commitment logic~\cite{Ommen2016}. However, MIQCP problems lack
the LP relaxation tightness of MILP, potentially increasing solve
times for large instances.

\subsection{Alternative approaches: data-driven surrogates}
\label{sec:rw_surrogates}

Recent work demonstrates that ANN-based surrogate models can replace
physics-based DH simulations with ${>}\,80\,\%$ time
savings~\cite{Guo2025}. However, surrogate approaches sacrifice
optimality guarantees, require retraining for each new network
configuration, and cannot provide the controlled, identical-instance
comparisons required to isolate modeling effects. For the purpose of
this study---where five formulations must be solved to global
optimality on identical problem instances---physics-based MILP/MIQCP
with commercial solvers remains the appropriate methodology.
Nevertheless, surrogates may complement the present framework in
operational settings: once the appropriate fidelity level has been
identified using the criteria derived here, a surrogate trained on
that level's outputs could enable real-time dispatch.

\subsection{Summary and positioning}
\label{sec:rw_summary}

\Cref{tab:prior_work} positions representative studies. The critical
observation is that no prior work isolates topology from physics, nor
quantifies linearization error against a quadratic reference in a
controlled DH dispatch setting.

\begin{table*}[!htbp]
\centering
\caption{Representative DH optimization studies. No prior study
isolates topology from physics or quantifies linearization error.}
\label{tab:prior_work}
\small
\setlength{\tabcolsep}{4pt}
\begin{tabular}{@{}l c c c c c c c@{}}
\toprule
Reference & Topology & Physics model & Formul.\ & COP &
  Scale & Temporal & Contr.?\ \\
\midrule
Connolly~\cite{Connolly2014_EnergyPLAN}
  & Copper.\ & Fixed~$\eta$ & LP & Const.\ & National & Annual & No \\
Gabrielli~\cite{Gabrielli2018}
  & Copper.\ & Fixed~$\eta$ & MILP & Const.\ & 10\,MW & Annual & No \\
Wirtz~\cite{Wirtz2020}
  & Zone\,(5) & Linear losses & MILP & Season.\ & 25\,MW & Annual & No \\
Mertz~\cite{Mertz2016}
  & Zone\,(8) & Lin.\ + $\Delta p$ & MINLP & Const.\
  & 20\,MW & Annual & No \\
Hari~\cite{Hari2024}
  & Full & DW + pumps & NLP & -- & 50\,MW & Steady & No \\
Guo~\cite{Guo2025}
  & Multi-sc.\ & Fixed~$\eta$ + ANN & MILP & Const.\ & 5--50\,MW
  & Annual & Partial$^{c}$ \\
Leitner~\cite{Leitner2019}
  & L1 + L3 & Differing phys.\ & MILP & Const.\ & 12\,MW
  & Annual & No \\
Bordin~\cite{Bordin2016}
  & Full & PWL~$\Delta p$ & MILP & -- & 40\,MW & Annual & No \\
Giraud~\cite{Giraud2015}
  & Full & Steady-state & NLP & $T$-dep.\ & 30\,MW & Annual & No \\
Vesterlund~\cite{Vesterlund2017}
  & Full & NLP thermal & NLP & $T$-dep.\ & 50\,MW & Annual & No \\
van der Heijde~\cite{vanderHeijde2017}
  & Full & Pipe model comp.\ & Various & -- & Review & --
  & Partial$^{a}$ \\
Blommaert~\cite{Blommaert2020}
  & Full & Adjoint NLP & NLP & $T$-dep.\ & 30\,MW & Design & No \\
\midrule
\textbf{This work}
  & \textbf{L1--$\Lnl$}
  & \textbf{Basic \& ext.}
  & \textbf{MILP+MIQCP}
  & \textbf{Pre-comp.}
  & \textbf{${\approx}\,12$\,MW}
  & \textbf{8760\,h}
  & \textbf{Yes} \\
\bottomrule
\end{tabular}

\vspace{2pt}
\raggedright\footnotesize
$^{a}$~Compared pipe loss formulations but not spatial aggregation
levels.\quad
$^{c}$~Compared spatial scales but not in a controlled same-network
experiment.
\end{table*}

Notably, Leitner et~al.~\cite{Leitner2019} compared copperplate and
full-graph representations for an Austrian DH network and found cost
differences of 3--8\,\%, but simultaneously changed the loss physics
between levels. None of these studies isolates either topology or
physics fidelity as the sole experimental variable.

\subsection{Thermal losses and storage dispatch}
\label{sec:rw_losses_storage}

Werner~\cite{Werner2017} reports European DH loss ratios of 5--25\,\%,
with modern pre-insulated networks typically below 10\,\%. Thermal
energy storage provides flexibility through peak shaving and load
shifting~\cite{Vandermeulen2018,Guelpa2019review}. Christidis et
al.~\cite{Christidis2017} demonstrate that optimal TES dispatch shifts
heat production by 4--8\,hours to exploit electricity price
differentials---a pattern directly dependent on accurate temporal cost
signals. In copperplate models, the only time-varying cost is the
electricity price. In network-aware models, losses introduce an
additional implicit holding cost: heat stored for later discharge must
additionally compensate for losses that are highest during cold
periods. Transport delays further modify storage dispatch by
introducing a temporal offset between injection and delivery.

\textbf{Hypothesis:} In networks where transport delay exceeds
$\SI{1}{\hour}$ or demand heterogeneity is high, storage dispatch
timing is a sensitive indicator of topology and physics fidelity
because temporal arbitrage depends on holding costs (invisible to L1)
and delivery timing (invisible to L3). For compact networks
($\tau_{\max} < \SI{1}{\hour}$, low~HI), timing differences are
expected to be small, confirming model equivalence rather than
insensitivity. This is tested in \Cref{sec:results_storage}.

%% ============================================================
%% 3. METHODOLOGY
%% ============================================================
\section{Methodology}
\label{sec:methodology}

This section presents the controlled experimental design
(\Cref{sec:exp_design}), the unified MILP formulation for L1/L2/L3
(\Cref{sec:milp}), the extended physics formulation for $\Lplus$
(\Cref{sec:milp_extended}), the quadratic reference formulation for
$\Lnl$ (\Cref{sec:miqcp}), the demand heterogeneity metric
(\Cref{sec:model_hi}), and implementation details
(\Cref{sec:model_impl}).

\subsection{Controlled experimental design}
\label{sec:exp_design}

\subsubsection{Three orthogonal comparison dimensions}

The experimental design isolates three effects through pairwise
comparisons (\Cref{fig:comparison_design}):
%
\begin{enumerate}[nosep]
  \item \textbf{Topology effect} (L1\,$\to$\,L2\,$\to$\,L3):
        Spatial resolution varies; basic physics held constant.
  \item \textbf{Physics fidelity effect} (L3\,$\to$\,$\Lplus$):
        Physics scope varies; topology and formulation class (MILP)
        held constant.
  \item \textbf{Linearization error}
        ($\Lplus$\,$\to$\,$\Lnl$): Mathematical formulation varies
        (MILP\,$\to$\,MIQCP); physics and topology held identical.
\end{enumerate}

\begin{figure}[htbp]
  \centering
  \includegraphics[width=\linewidth]{F1_experimental_design.pdf}
  \caption{Experimental design with three orthogonal comparison
  dimensions. Each arrow isolates exactly one modeling dimension.}
  \label{fig:comparison_design}
\end{figure}

\subsubsection{Basic physics (L1/L2/L3)}

The basic physics model is held \emph{identical} across L2 and L3:
%
\begin{itemize}[nosep]
  \item \textbf{Steady-state thermal losses:} per-pipe, separately
        for supply and return.
  \item \textbf{Fixed supply temperature:} determined by a heating
        curve, not optimized.
  \item \textbf{Pre-computed COP:} time-varying parameters from
        analytical Carnot-based model (\Cref{app:cop}).
  \item \textbf{No transport delays:} instantaneous heat propagation.
  \item \textbf{No pressure optimization:} pumping power excluded.
  \item \textbf{Dispatch-only:} all component capacities fixed.
\end{itemize}

\textbf{Rationale for fixed supply temperature:} Treating
$T_{\mathrm{sup}}$ as a decision variable would introduce bilinear
$T \cdot \dot{m}$ terms in the basic loss model, destroying MILP
linearity for L2/L3. The pre-determined heating curve is standard
industrial practice for the case study network and ensures
comparability across levels.

\subsubsection{Extended physics ($\Lplus$\,/\,$\Lnl$)}

The extended physics model adds three phenomena to the L3 base:
%
\begin{itemize}[nosep]
  \item \textbf{Pressure drop:} Darcy--Weisbach friction losses and
        pumping power (\Cref{sec:model_pressure}).
  \item \textbf{Temperature-dependent loss propagation:} downstream
        node temperatures depend on upstream pipe losses
        (\Cref{sec:model_temp_prop}).
  \item \textbf{Transport delay:} finite heat propagation time
        through pipes (\Cref{sec:model_delay}).
\end{itemize}

In $\Lplus$, these phenomena are linearized for MILP compatibility.
In $\Lnl$, they are modeled with native quadratic/bilinear
constraints (MIQCP), serving as the highest-fidelity tractable
reference within the MIQCP class.

\subsubsection{Model scope}
\label{sec:model_scope}

\Cref{tab:model_scope} delineates which physical phenomena are
included in each formulation level.

\begin{table*}[!htbp]
\centering
\caption{Model scope across all five variants. ``PWL'' indicates
piecewise-linear approximation; ``Quad/Bilin'' indicates native
quadratic or bilinear constraint.}
\label{tab:model_scope}
\begin{tabular}{@{} l C{1.0cm} C{1.0cm} C{1.0cm} C{1.0cm} C{1.0cm} @{}}
\toprule
Physical phenomenon & L1 & L2 & L3 & $\Lplus$ & $\Lnl$ \\
\midrule
\multicolumn{6}{@{}l}{\textit{Thermal}} \\[2pt]
\quad Steady-state pipe losses ($UL\Delta T$)
  & -- & \checkmark & \checkmark & \checkmark & \checkmark \\
\quad Temperature propagation (downstream $T$ reduction)
  & -- & -- & -- & PWL & Bilin \\
\quad Time-varying COP (pre-computed)
  & \checkmark & \checkmark & \checkmark & \checkmark & \checkmark \\
\quad Storage self-discharge \& efficiency
  & \checkmark & \checkmark & \checkmark & \checkmark & \checkmark \\
\midrule
\multicolumn{6}{@{}l}{\textit{Hydraulic}} \\[2pt]
\quad Pressure drop (Darcy--Weisbach)
  & -- & -- & -- & PWL & Quad \\
\quad Pumping power (lumped)
  & -- & -- & -- & \checkmark & \checkmark \\
\midrule
\multicolumn{6}{@{}l}{\textit{Temporal}} \\[2pt]
\quad Transport delay (plug-flow, fixed)
  & -- & -- & -- & \checkmark & \checkmark \\
\quad Hourly grid emission factors
  & \checkmark & \checkmark & \checkmark & \checkmark & \checkmark \\
\midrule
\multicolumn{6}{@{}l}{\textit{Excluded across all levels}} \\[2pt]
\quad Sub-hourly transient thermal inertia
  & \multicolumn{5}{c}{--} \\
\quad Investment / capacity sizing
  & \multicolumn{5}{c}{--} \\
\quad Supply temperature optimization
  & \multicolumn{5}{c}{--} \\
\quad Bidirectional / ring flow
  & \multicolumn{5}{c}{--} \\
\bottomrule
\end{tabular}
\end{table*}

\subsubsection{Level definitions}
\label{sec:topology_def}

\paragraph{Level~1 (Copperplate)}
Single bus, no spatial structure, no pipe losses
($Q_{\mathrm{loss}}=0$ for all~$t$). All assets co-located at a
single node. A single scalar heat balance.

\paragraph{Level~2 (Simplified multi-node)}
Aggregated demand zones (7~zones in the primary case) connected by
simplified pipe segments preserving the tree topology. The total
$U \cdot L$ product is preserved from L3, ensuring matched aggregate
annual loss energy. Asset locations are preserved at their physical
positions (\Cref{sec:case_l2_zones}).

\paragraph{Level~3 (Detailed multi-node, basic MILP)}
Full spatial resolution (15~junction nodes in the primary case).
Individual pipe parameters. Same basic loss formulation as L2. No
pressure drop, no temperature propagation, no transport delay.

\paragraph{$\Lplus$ (Detailed multi-node, extended MILP)}
Same topology as L3. Adds linearized pressure drop~(PWL),
temperature propagation (Taylor linearization), and transport delay
(fixed). All nonlinear terms are approximated to preserve MILP
solvability.

\paragraph{$\Lnl$ (Detailed multi-node, extended MIQCP)}
Same topology and physics scope as $\Lplus$. Pressure drop enters as
a native quadratic constraint; temperature propagation involves
bilinear products handled by the MIQCP solver. Both $\Lplus$ and
$\Lnl$ share the PWL approximation of the exponential decay factor
(not representable in MIQCP); the linearization error quantified is
therefore limited to the bilinear temperature-mixing and quadratic
pressure-drop terms.

\subsection{Basic MILP formulation (L1/L2/L3)}
\label{sec:milp}

\subsubsection{Objective function}

Total annual operational cost is minimized:
%
\begin{equation}
  \min \; Z = C^{\mathrm{en}} + C^{\mathrm{fuel}}
    + C^{\mathrm{dump}} + C^{\ce{CO2}} + C^{\mathrm{dem}}
    + C^{\mathrm{pump}}
  \label{eq:objective}
\end{equation}
%
where $C^{\mathrm{pump}} = 0$ for L1/L2/L3 and is defined in
\Cref{sec:model_pressure} for $\Lplus$/$\Lnl$.

\textbf{Electricity procurement and sales:}
\begin{equation}
  C^{\mathrm{en}} = \Delta t \sum_{t \in \mathcal{T}} \left(
    P_t^{\mathrm{buy}} \lambda_t^{\mathrm{buy}}
    - P_t^{\mathrm{sell}} \lambda_t^{\mathrm{sell}} \right)
  \label{eq:cost_energy}
\end{equation}

\textbf{Fuel cost:}
\begin{equation}
  C^{\mathrm{fuel}} = \Delta t
    \sum_{g \in \mathcal{G}^{\mathrm{fuel}}}
    \sum_{t \in \mathcal{T}} F_{g,t}\, \lambda_g^{\mathrm{fuel}}
  \label{eq:cost_fuel}
\end{equation}

\textbf{Heat curtailment penalty:}
\begin{equation}
  C^{\mathrm{dump}} = \Delta t \sum_{n \in \mathcal{N}}
    \sum_{t \in \mathcal{T}} Q_{n,t}^{\mathrm{dump}}\,
    \lambda^{\mathrm{dump}}
  \label{eq:cost_dump}
\end{equation}

\textbf{\ce{CO2} cost:}
\begin{equation}
  C^{\ce{CO2}} = \frac{\lambda^{\ce{CO2}}}{1000} \left(
    \sum_{t} E_t^{\ce{CO2},\mathrm{grid}}
    + \sum_{g,t} E_{g,t}^{\ce{CO2},\mathrm{fuel}} \right)
  \label{eq:cost_co2}
\end{equation}

\textbf{Demand charge} (annual peak import fee):
\begin{equation}
  C^{\mathrm{dem}} = \lambda^{\mathrm{dem}} \cdot f^{\mathrm{yr}}
    \cdot P^{\mathrm{peak}}
  \label{eq:cost_demand}
\end{equation}

\subsubsection{Energy balance constraints}
\label{sec:model_balance}

\textbf{L1~--- Copperplate balance (no losses):}
\begin{equation}
  \sum_{g} Q_{g,t}^{\mathrm{th}}
    + \sum_{s} Q_{s,t}^{\mathrm{dis}}
  = Q_t^{\mathrm{dem}} + Q_t^{\mathrm{dump}}
    + \sum_{s} Q_{s,t}^{\mathrm{ch}}
  \qquad \forall\, t
  \label{eq:balance_L1}
\end{equation}

\textbf{L2/L3/$\Lplus$/$\Lnl$~--- Nodal balance:}
\begin{multline}
  \sum_{p \in \mathcal{P}_n^{\mathrm{in}}}
    \bigl(Q_{p,t} - Q_{p,t}^{\mathrm{loss,pipe}}\bigr)
  + \sum_{g \in \mathcal{G}_n} Q_{g,t}^{\mathrm{th}}
  + \sum_{s \in \mathcal{S}_n} Q_{s,t}^{\mathrm{dis}} \\
  = Q_{n,t}^{\mathrm{dem}} + Q_{n,t}^{\mathrm{dump}}
  + \sum_{s \in \mathcal{S}_n} Q_{s,t}^{\mathrm{ch}}
  + \sum_{p \in \mathcal{P}_n^{\mathrm{out}}} Q_{p,t}
  \quad \forall\, n \in \mathcal{N},\; \forall\, t
  \label{eq:balance_nodal}
\end{multline}

In the primary case, generation assets are distributed: CHP, storage,
gas boiler, and biomass boiler at~j$_1$; heat pump and electrode
boiler at~j$_{12}$. For nodes without local generation, the
corresponding terms $\mathcal{G}_n$ and $\mathcal{S}_n$ are empty.
The contrast between Eqs.~\eqref{eq:balance_L1} and
\eqref{eq:balance_nodal} encodes the central experimental variable:
L1's missing loss and spatial flow terms are the structural
differences driving the L1\,$\to$\,L2 gap.

\textbf{Global electricity balance:}
\begin{equation}
  P_t^{\mathrm{buy}}
    + \sum_{g \in \mathcal{G}^{\mathrm{CHP}}} P_{g,t}^{\mathrm{el}}
  = \sum_{h \in \mathcal{H}} P_{h,t}^{\mathrm{el}}
    + \sum_{r \in \mathcal{R}} P_{r,t}^{\mathrm{el}}
    + P_t^{\mathrm{pump}}
    + P_t^{\mathrm{sell}}
  \qquad \forall\, t
  \label{eq:balance_el}
\end{equation}

\textbf{Grid import/export mutual exclusivity:}
\begin{equation}
  P_t^{\mathrm{buy}} \leq M^{\mathrm{grid}}\, z_t, \qquad
  P_t^{\mathrm{sell}} \leq M^{\mathrm{grid}}\, (1 - z_t)
  \qquad \forall\, t
  \label{eq:grid_bigm}
\end{equation}

\subsubsection{Thermal loss model --- basic (L2/L3)}
\label{sec:model_losses}

Losses are computed separately for supply and return pipes using
steady-state conduction with \emph{fixed} temperatures:
%
\begin{align}
  Q_{p,t}^{\mathrm{loss,sup}} &= \frac{U_p L_p \bigl(
    T_{\mathrm{sup}}^{\mathrm{nom}}(t) - T_{\mathrm{ground}}(t)
    \bigr)}{10^6}
  \quad [\text{MW}]
  \label{eq:loss_supply} \\[4pt]
  Q_{p,t}^{\mathrm{loss,ret}} &= \frac{U_p L_p \bigl(
    T_{\mathrm{ret}}^{\mathrm{nom}} - T_{\mathrm{ground}}(t)
    \bigr)}{10^6}
  \quad [\text{MW}]
  \label{eq:loss_return}
\end{align}
%
where $U_p$\,[\si{\watt\per\metre\per\kelvin}] is the linear heat
transfer coefficient, $L_p$\,[m] the pipe length, and
$T_{\mathrm{ground}}(t)$ the ground temperature. These losses enter
as \emph{fixed parameters}, not decision variables, because
temperatures are pre-determined. This is the key simplification
preserving MILP linearity. For L1,
$Q_{p,t}^{\mathrm{loss}} = 0$.

\subsubsection{Component models}
\label{sec:model_components}

\paragraph{Heat pumps}
Each heat pump~$h$ supports waste-heat recovery~(WRG) with
time-varying COP and a default source with constant COP:
%
\begin{align}
  Q_{h,t} &= Q_{h,t}^{\mathrm{wrg}} + Q_{h,t}^{\mathrm{def}}
  \label{eq:hp_split} \\
  P_{h,t}^{\mathrm{el}} &= \frac{Q_{h,t}^{\mathrm{wrg}}}{
    \mathrm{COP}_{h,t}^{\mathrm{wrg}}}
    + \frac{Q_{h,t}^{\mathrm{def}}}{\mathrm{COP}_h^{\mathrm{def}}}
  \label{eq:hp_pel}
\end{align}
%
Capacity and minimum part-load:
\begin{equation}
  \phi_h^{\mathrm{min}}\, \bar{Q}_h\, u_{h,t}
  \leq Q_{h,t}
  \leq \bar{Q}_h\, u_{h,t}
  \qquad \forall\, h,\; \forall\, t
  \label{eq:hp_cap}
\end{equation}

\paragraph{COP treatment}
In a full optimization where source and sink temperatures are decision
variables, heat pump output couples nonlinearly. Three linearization
strategies exist: (1)~constant COP (up to 30\,\%
error~\cite{Ommen2016}); (2)~McCormick envelopes; (3)~pre-computed
COP time series. This work adopts strategy~(3):
$\mathrm{COP}_{h,t}^{\mathrm{wrg}}$ is evaluated offline from
exogenous source temperature data using an analytical Lorenz-cycle
model (\Cref{app:cop}), entering the MILP as a fixed parameter with
a combined error of ${\approx}\,3.6\,\%$---well below all topology
and physics gaps.

\paragraph{Thermal storage}
State-of-charge dynamics:
\begin{equation}
  E_{s,t} = (1-\alpha_s)^{\Delta t}\, E_{s,t-1}
    + \eta_{s}^{\mathrm{ch}}\, Q_{s,t}^{\mathrm{ch}}\, \Delta t
    - \frac{Q_{s,t}^{\mathrm{dis}}\, \Delta t}{
      \eta_{s}^{\mathrm{dis}}}
  \qquad \forall\, s,\; \forall\, t
  \label{eq:storage_soc}
\end{equation}

\paragraph{Thermal generators and CHP}
Standard linear efficiency constraints (full formulation in
\Cref{app:component_models}). The CHP is modeled as a backpressure
unit with fixed heat-to-power ratio.

\subsubsection{Emissions accounting}

\begin{align}
  E_t^{\ce{CO2},\mathrm{grid}} &=
    P_t^{\mathrm{buy}} \cdot e_t^{\mathrm{grid}} \cdot \Delta t
  \label{eq:em_grid} \\
  E_{g,t}^{\ce{CO2},\mathrm{fuel}} &=
    F_{g,t} \cdot e_g^{\mathrm{fuel}} \cdot \Delta t
  \label{eq:em_fuel}
\end{align}

%% ============================================================
%% 3.3 EXTENDED PHYSICS (L3+ MILP)
%% ============================================================
\subsection{Extended physics formulation --- $\Lplus$ (MILP)}
\label{sec:milp_extended}

$\Lplus$ augments L3 with three physical phenomena, each linearized
to preserve MILP solvability.

\subsubsection{Pressure drop model}
\label{sec:model_pressure}

The Darcy--Weisbach equation:
%
\begin{equation}
  \Delta p_{p} = f_D \,\frac{L_p}{D_p}\,
    \frac{8\,\dot{m}_p^2}{\rho^2 \pi^2 D_p^4}
  \qquad [\text{Pa}]
  \label{eq:darcy}
\end{equation}

Mass flow is linked to thermal power:
\begin{equation}
  \dot{m}_{p,t} = \frac{Q_{p,t}}{c_p \bigl(
    T_{\mathrm{sup}}^{\mathrm{nom}}(t)
    - T_{\mathrm{ret}}^{\mathrm{nom}}(t)\bigr)}
  \label{eq:massflow}
\end{equation}

Since temperatures are pre-computed parameters, $\dot{m}_{p,t}$ is
linear in~$Q_{p,t}$, but pressure drop is \emph{quadratic}
in~$Q_{p,t}$.

\paragraph{MILP linearization (PWL)}
The quadratic pressure drop is approximated using a piecewise-linear
function with $K=3$ breakpoints:
\begin{equation}
  \Delta p_{p,t}^{\mathrm{PWL}} = \sum_{k=1}^{K}
    m_k^{(p)} \cdot w_{p,t,k}
  \label{eq:pressure_pwl}
\end{equation}
subject to SOS2 constraints on~$w_{p,t,k}$. The worst-case
approximation error is bounded by
$\varepsilon_{\Delta p}^{\max} = R_p \bar{Q}_p^2 / (4K^2)$, i.e.,
2.8\,\% of peak pressure drop for $K=3$ (derivation in
\Cref{app:pwl}).

\paragraph{Choice of $K=3$}
The number of PWL segments balances approximation accuracy against
variable count. With $K=3$, the worst-case pressure-drop error is
2.8\,\% of peak. A sensitivity analysis on $K \in \{3, 5, 8\}$
(\Cref{app:pwl}) confirms diminishing returns: $K=5$ achieves 1\,\%
accuracy, but the linearization error in total-cost terms remains
$+0.35$\,\% at $K=3$, which is within the accepted tolerance for
planning-level decisions.

\paragraph{Total network pumping power}
\begin{equation}
  P_t^{\mathrm{pump}} = \frac{1}{\eta_{\mathrm{pump}}}
    \sum_{p \in \mathcal{P}}
    \frac{\dot{m}_{p,t} \cdot \Delta p_{p,t}}{\rho}
  \label{eq:pump_power}
\end{equation}

In $\Lplus$, pumping power is directly piecewise-linearized as a
function of~$Q_{p,t}$ (cubic relationship) with $K=3$ breakpoints:
\begin{equation}
  P_{p,t}^{\mathrm{pump,PWL}} = \sum_{k=1}^{K}
    c_k^{(p)} \cdot w_{p,t,k}
  \label{eq:pump_pwl}
\end{equation}

\textbf{Lumped pumping model:} Pumping power is computed as a network
aggregate rather than assigned to specific pump locations. For
networks with distributed pumps, this may overestimate global
efficiency~\cite{Hari2024}. For the case study's centrally-pumped tree
network, the lumped model is exact.

\subsubsection{Temperature-dependent loss propagation}
\label{sec:model_temp_prop}

The steady-state temperature at the outlet of pipe~$p$ is:
%
\begin{equation}
  T_{p,t}^{\mathrm{out}} = T_{\mathrm{ground}}(t)
    + \bigl(T_{n_1,t}^{\mathrm{sup}} - T_{\mathrm{ground}}(t)\bigr)
    \cdot \underbrace{\exp\!\left(
      \frac{-U_p L_p}{\dot{m}_{p,t}\, c_p}
    \right)}_{\displaystyle\phi_p(\dot{m}_{p,t})}
  \label{eq:temp_prop}
\end{equation}

This introduces two nonlinearities: (1)~the exponential decay
$\phi_p(\dot{m})$ and (2)~the bilinear product
$T_{n_1,t}^{\mathrm{sup}} \cdot \phi_{p,t}$.

\paragraph{MILP linearization}
\emph{Step~1~--- PWL decay factor:}
\begin{equation}
  \phi_{p,t}^{\mathrm{PWL}} = \sum_{k=1}^{K}
    \phi_p(\dot{m}_k) \cdot w_{p,t,k}
  \label{eq:decay_pwl}
\end{equation}

\emph{Step~2~--- Taylor linearization of bilinear product:}
\begin{equation}
  T_{p,t}^{\mathrm{out}} \approx T_{\mathrm{ground}}(t)
    + \bigl(T^{\mathrm{nom}}(t) - T_{\mathrm{ground}}(t)\bigr)
      \cdot \phi_{p,t}^{\mathrm{PWL}}
    + \phi^{\mathrm{nom}}_p
      \cdot \bigl(T_{n_1,t}^{\mathrm{sup}} - T^{\mathrm{nom}}(t)\bigr)
  \label{eq:temp_prop_lin}
\end{equation}

The extended pipe loss:
\begin{equation}
  Q_{p,t}^{\mathrm{loss,ext}} \approx
    \frac{U_p L_p}{10^6} \cdot \bar{T}_{p,t}^{\Delta}
  \label{eq:loss_ext_lin}
\end{equation}
where $\bar{T}_{p,t}^{\Delta}$ is the effective mean temperature
excess at linearized temperatures (derivation in
\Cref{app:taylor_derivation}).

\subsubsection{Transport delay model}
\label{sec:model_delay}

Heat injected at node~$n_1$ at time~$t$ arrives at the downstream end
of pipe~$p$ at time $t + \tau_p$:
%
\begin{equation}
  \tau_p = \frac{L_p}{v_p^{\mathrm{nom}}}
         = \frac{L_p\, \rho\, A_p}{\dot{m}_p^{\mathrm{nom}}}
  \label{eq:delay}
\end{equation}

Discretized to integer timesteps:
$k_p = \mathrm{round}(\tau_p / \Delta t)$. The delayed energy
balance:
\begin{equation}
  Q_{p,t}^{\mathrm{delivered}} =
    Q_{p,\,t-k_p} - Q_{p,\,t-k_p}^{\mathrm{loss,pipe}}
  \qquad \forall\, p,\; \forall\, t \geq k_p + 1
  \label{eq:delayed_balance}
\end{equation}

\textbf{Validity range:} The fixed-delay approximation is most
accurate when flow variations are moderate ($\pm 30\,\%$ of nominal).
For the primary case's longest path (\SI{2125}{\metre} at
${\sim}\SI{1.0}{\metre\per\second}$), $\tau \approx 0.6$\,h,
discretized to $k_p = 1$. This is marginal for the primary case but
becomes significant for the synthetic analysis where $\tau_{\max}$
reaches \SI{5.2}{\hour}.

%% ============================================================
%% 3.4 QUADRATIC REFERENCE (L3^NL MIQCP)
%% ============================================================
\subsection{Quadratic reference formulation --- $\Lnl$ (MIQCP)}
\label{sec:miqcp}

$\Lnl$ uses the same topology, component models, and boundary
conditions as $\Lplus$, but replaces piecewise-linear and first-order
approximations with native quadratic and bilinear constraints.

\textbf{Scope of $\Lnl$ as reference:} $\Lnl$ is not a physical
ground truth---both $\Lplus$ and $\Lnl$ share the PWL approximation
of the exponential decay factor. The linearization error quantified is
limited to: (1)~quadratic vs.\ PWL pressure drop, and (2)~bilinear
vs.\ Taylor-linearized temperature mixing.

\paragraph{Pressure drop --- quadratic}
\begin{equation}
  \Delta p_{p,t} = R_p \cdot Q_{p,t}^2
  \qquad \forall\, p,\; \forall\, t
  \label{eq:darcy_quad}
\end{equation}

\paragraph{Pumping power}
Via auxiliary variable $W_{p,t} = Q_{p,t}^2$:
\begin{equation}
  P_{p,t}^{\mathrm{pump}} =
    \frac{R_p}{\eta_{\mathrm{pump}}\, \rho\, c_p\,
      \Delta T_{\mathrm{nom}}(t)}
    \cdot Q_{p,t} \cdot W_{p,t}
  \label{eq:pump_bilinear}
\end{equation}

\paragraph{Temperature propagation --- bilinear}
\begin{equation}
  \Theta_{p,t} = T_{n_1,t}^{\mathrm{sup}}
    \cdot \phi_{p,t}^{\mathrm{PWL}}
  \label{eq:bilinear_temp}
\end{equation}

\paragraph{Transport delay}
Identical to $\Lplus$ (fixed).

\Cref{tab:lplus_vs_lnl} summarizes the mathematical treatment across
the two extended formulations.

\begin{table}[htbp]
\centering
\caption{Mathematical treatment of extended physics.}
\label{tab:lplus_vs_lnl}
\small
\begin{tabular}{@{} l l l @{}}
\toprule
Phenomenon & $\Lplus$ (MILP) & $\Lnl$ (MIQCP) \\
\midrule
Pressure drop   & PWL ($K{=}3$)       & Quadratic $RQ^2$ \\
Pumping power   & PWL of cubic        & Bilinear $Q \cdot W$ \\
Temp.\ decay~$\phi$ & PWL of exponential & PWL of exp.$^{*}$ \\
Temp.\ mixing   & Taylor linearization & Bilinear $T \cdot \phi$ \\
Transport delay & Fixed               & Fixed \\
\midrule
Optimality      & Global (LP relax.)  & Global (spatial branch.) \\
\bottomrule
\end{tabular}

\vspace{2pt}
\raggedright\footnotesize
$^{*}$Shared approximation; not part of measured linearization error.
\end{table}

\subsection{Demand heterogeneity index}
\label{sec:model_hi}

\begin{equation}
  d_n = \frac{\sum_t Q_{n,t}^{\mathrm{dem}}}{\sum_n \sum_t
    Q_{n,t}^{\mathrm{dem}}},
  \qquad
  \mathrm{HI} = \frac{N}{N-1} \sum_{n=1}^{N}
    \left(d_n - \frac{1}{N}\right)^{\!2}
  \label{eq:hi}
\end{equation}

$\mathrm{HI} = 0$ indicates uniform demand;
$\mathrm{HI} = 1$ indicates complete concentration.

\subsection{Implementation}
\label{sec:model_impl}

The model is implemented in Python~3.10 using Pyomo~\cite{pyomo}.
All five variants are instantiated from a single parameterized
codebase via YAML configuration files. Solver: Gurobi~10.0, with MIP
gap ${\leq}\,0.5\,\%$ for MILP and ${\leq}\,0.5\,\%$ for MIQCP
(\texttt{NonConvex=2}). Identical random seeds.
L1/L2/L3/$\Lplus$ are solved from scratch independently.
$\Lnl$ is warm-started from the $\Lplus$ solution to improve
MIQCP branch-and-bound convergence; the reported linearization
error is based on gap-certified terminal costs, not the
warm-start point, so the comparison remains unbiased.

\textbf{Problem scaling.} The $\Lplus$ formulation adds $K$~SOS2
weights, 1~auxiliary flow variable, and 1~node temperature variable
per pipe per timestep:
$|\mathcal{P}| \times |\mathcal{T}| \times (K+2)
= 14 \times 8{,}760 \times 10 \approx 1{,}226{,}000$
additional continuous variables. The $\Lnl$ formulation adds
$2 \times 14 \times 8{,}760 \approx 245{,}000$ quadratic constraints.

%% ============================================================
%% 4. CASE STUDIES
%% ============================================================
\section{Case studies}
\label{sec:case_studies}

\subsection{Primary case: industrial district heating network}
\label{sec:case_stadtbach}

\subsubsection{Network topology and components}

The primary case is an industrial DH network operated by a
medium-sized utility in southern Germany. The network has a radial
tree structure with a central production node~(j$_1$) and two
distribution arms (15~junctions, 14~pipe segments, 27~consumers). The
heat pump and electrode boiler are located at junction~j$_{12}$ on the
East arm, approximately \SI{1680}{\metre} from the production node.
This decentralized placement is operationally motivated: j$_{12}$ is
adjacent to the industrial waste-heat source supplying the HP
evaporator. It introduces a structural asymmetry that amplifies the
topology effect: heat produced at~j$_{12}$ does not traverse the
trunk, reducing losses relative to centralized production---an
advantage invisible to L1.

\Cref{tab:case_summary} consolidates key parameters. Absolute demand
values are anonymized per non-disclosure agreement; all results are
reported as relative differences between model levels.

\begin{table}[htbp]
\centering
\caption{Primary case study: consolidated parameters.}
\label{tab:case_summary}
\small
\begin{tabular}{@{} l l l @{}}
\toprule
Parameter & Value & Notes \\
\midrule
\multicolumn{3}{@{}l}{\textit{Network topology}} \\[2pt]
Junction nodes / consumers   & 15\,/\,27 & Tree structure \\
Pipe segments                & 14 & \\
Total pipe length (one-way)  & \SI{3255}{\metre} & \\
Longest path (j$_1 \to$ j$_{15}$) & \SI{2125}{\metre} & \\
Max.\ transport delay ($\tau$)
  & ${\approx}\,\SI{35}{\minute}$ & $k_p = 1$ \\[4pt]
\midrule
\multicolumn{3}{@{}l}{\textit{Thermal parameters}} \\[2pt]
Supply temperature
  & 70--\SI{100}{\celsius} & PWL heating curve \\
Return temperature
  & 40--\SI{55}{\celsius} & \\
Ground temperature
  & \SI{10}{\celsius} & Fixed \\
$U$-value range (sup\,/\,ret)
  & 0.28--0.32\,/\,0.30--0.34\,\si{\watt\per\metre\per\kelvin}
  & Per-pipe \\[4pt]
\midrule
\multicolumn{3}{@{}l}{\textit{Hydraulic ($\Lplus$/$\Lnl$ only)}} \\[2pt]
Pipe roughness     & \SI{0.5}{\milli\metre} & Aged steel \\
Pump efficiency    & 0.75 & Central station \\
PWL breakpoints    & $K = 3$ & \\[4pt]
\midrule
\multicolumn{3}{@{}l}{\textit{Generation assets at~j$_1$}} \\[2pt]
CHP thermal\,/\,electrical
  & 0.2\,/\,\SI{0.1}{\mega\watt} & Backpressure$^{\ddagger}$ \\
Gas boiler
  & \SI{13}{\mega\watt_{\mathrm{th}}} & $\eta = 0.90$ \\
Biomass boiler
  & \SI{3.3}{\mega\watt_{\mathrm{th}}} & $\eta = 0.85$ \\
Thermal storage
  & \SI{500}{\mega\watt\hour}\,/\,\SI{30}{\mega\watt} & Cyclic BC \\[4pt]
\midrule
\multicolumn{3}{@{}l}{\textit{Generation assets at~j$_{12}$ (East arm)}} \\[2pt]
Heat pump (WRG)
  & \SI{5}{\mega\watt_{\mathrm{th}}} & Lorenz $\eta = 0.6$ \\
Electrode boiler
  & \SI{5}{\mega\watt_{\mathrm{th}}} & $\eta = 0.95$ \\[4pt]
\midrule
\multicolumn{3}{@{}l}{\textit{Demand}} \\[2pt]
Peak demand
  & ${\approx}\,\SI{12}{\mega\watt}$ & \\
Demand heterogeneity (HI)
  & $0.05$ & \\
\bottomrule
\end{tabular}
\vspace{2pt}\raggedright\footnotesize
$^{\ddagger}$CHP supplies ${<}\,2\,\%$ of annual thermal output
relative to the 12\,MW peak. It is modeled for completeness but does
not materially influence any of the level-comparison results.
\end{table}

\begin{figure}[htbp]
  \centering
  \includegraphics[width=\linewidth]{F2_network_topology.pdf}
  \caption{Network topology. Assets at j$_1$ (CHP, TES, boilers)
  and j$_{12}$ (HP, electrode boiler). Transport delays annotated.}
  \label{fig:topology}
\end{figure}

\subsubsection{Load-dependent temperature profile}
\label{sec:case_temp_profile}

Supply and return temperatures follow a piecewise-linear heating curve
(\Cref{tab:pwl}). In L1/L2/L3, these define loss parameters. In
$\Lplus$/$\Lnl$, they additionally define the nominal profile for
temperature propagation.

\begin{table}[htbp]
\centering
\caption{PWL temperature profile (heating curve).}
\label{tab:pwl}
\begin{tabular}{@{} l S[table-format=1.1] S[table-format=1.1]
                     S[table-format=1.1] S[table-format=1.1] @{}}
\toprule
Load fraction & {0.0} & {0.3} & {0.6} & {1.0} \\
\midrule
$T_{\mathrm{sup}}$ [\si{\celsius}] & 70 & 80 & 90 & 100 \\
$T_{\mathrm{ret}}$ [\si{\celsius}] & 55 & 50 & 45 & 40 \\
$\Delta T$ [\si{\celsius}]         & 15 & 30 & 45 & 60 \\
\bottomrule
\end{tabular}
\end{table}

\subsubsection{L2 zone construction}
\label{sec:case_l2_zones}

L2 aggregates the 15~L3 junctions into 7~zones by merging adjacent
nodes that share a branch segment, preserving the tree topology
(\Cref{tab:l2_zones}). Three controlled-matching principles govern
the construction:
%
\begin{enumerate}[nosep]
  \item \textbf{$\sum U_p L_p$ preservation:} Each L2 pipe carries
        the exact sum of the $U \cdot L$ products of the L3 pipes it
        replaces (total: \SI{1011}{\watt\per\kelvin}, matched to
        ${<}\,0.01\,\%$).
  \item \textbf{Asset-location preservation:} HP and electrode boiler
        remain at Zone~F (corresponding to~j$_{12}$), not relocated
        to~j$_1$.
  \item \textbf{Individual demand profiles:} Each zone sums the
        individual consumer time series of its constituent nodes,
        preserving temporal heterogeneity.
\end{enumerate}

\begin{table}[htbp]
\centering
\caption{L2 zone mapping and controlled matching to L3.}
\label{tab:l2_zones}
\small
\begin{tabular}{@{} l l l l l @{}}
\toprule
Zone & L3 nodes & Consumers & Assets & Pipe to [m] \\
\midrule
A & j$_1$          & V$_1$              & CHP, TES, Gas, Bio & -- \\
B & j$_2$, j$_3$   & V$_2$--V$_3$       & --       & A$\to$B: 650 \\
C & j$_4$, j$_5$   & V$_4$--V$_9$       & --       & B$\to$C: 500 \\
D & j$_6$--j$_8$   & V$_{10}$--V$_{14}$ & --       & C$\to$D: 450 \\
E & j$_9$--j$_{11}$& V$_{15}$--V$_{18}$ & --       & B$\to$E: 880 \\
F & j$_{12}$, j$_{13}$ & V$_{19}$--V$_{24}$ & HP, EBoiler
  & E$\to$F: 470 \\
G & j$_{14}$, j$_{15}$ & V$_{25}$--V$_{27}$ & --
  & F$\to$G: 305 \\
\bottomrule
\end{tabular}
\end{table}

\subsubsection{Controlled matching across all five variants}
\label{sec:case_matching}

\Cref{tab:matching} documents the controlled matching explicitly.

\begin{table*}[!htbp]
\centering
\caption{Controlled matching. Topology is the sole variable for
L1\,$\to$\,L2\,$\to$\,L3; physics scope for L3\,$\to$\,$\Lplus$;
formulation for $\Lplus$\,$\to$\,$\Lnl$.}
\label{tab:matching}
\small
\begin{tabular}{@{} l C{1.4cm} C{1.4cm} C{1.4cm} C{1.8cm} C{1.8cm} @{}}
\toprule
Parameter & L1 & L2 (7z) & L3 & $\Lplus$ & $\Lnl$ \\
\midrule
Nodes                   & 1   & 7    & 15   & 15         & 15 \\
Pipe segments           & 0   & 6    & 14   & 14         & 14 \\
Topology structure      & --  & Tree & Tree & Tree       & Tree \\
Thermal losses          & None & Physics & Physics & Physics & Physics \\
Pressure drop           & --  & --   & --   & PWL ($K{=}3$) & Quadratic \\
Temp.\ propagation      & --  & --   & --   & Linearized & Bilinear \\
Transport delay         & --  & --   & --   & Fixed      & Fixed \\
$\sum U_p L_p$ [W/K]   & 0   & 1011 & 1011 & 1011       & 1011 \\
HP / EBoiler location   & j$_1^{*}$ & Zone~F & j$_{12}$ & j$_{12}$ & j$_{12}$ \\
CHP/TES/GB/Bio location & j$_1$ & Zone~A & j$_1$ & j$_1$  & j$_1$ \\
Total demand per~$t$    & $=$ L3 & $=$ L3 & (ref.) & $=$ L3 & $=$ L3 \\
Demand representation   & Aggreg. & Zone sums & Per-node & Per-node & Per-node \\
COP time series         & $=$ L3 & $=$ L3 & (ref.) & $=$ L3 & $=$ L3 \\
Prices \& emission fact.\ & $=$ L3 & $=$ L3 & (ref.) & $=$ L3 & $=$ L3 \\
Temporal resolution     & 1\,h & 1\,h & 1\,h & 1\,h & 1\,h \\
\midrule
\textbf{Formulation}    & MILP & MILP & MILP & MILP & MIQCP \\
Solver / MIP gap
  & \multicolumn{4}{c}{Gurobi~10.0\,/\,0.5\,\%}
  & Gurobi\,/\,0.5\,\% \\
\bottomrule
\end{tabular}

\vspace{2pt}
\raggedright\footnotesize
$^{*}$L1 has only 1~node; all assets are co-located by construction.
The HP location effect is part of what L1\,$\to$\,L2 measures.
\end{table*}

\subsection{Validation protocol}
\label{sec:case_validation}

\subsubsection{Rationale for split validation}

The network underwent a technology upgrade (HP, electrode boiler, TES
installation) \emph{after} the measurement record was collected.
Following Ku\'{s} et~al.~\cite{Kus2025}, we adopt a two-stage
strategy:
%
\begin{enumerate}[nosep]
  \item \textbf{Stage~1~--- Direct network validation:} Pipe model
        validated against pre-upgrade monitoring data; new assets set
        to zero capacity.
  \item \textbf{Stage~2~--- Necessary-condition verification:}
        Dispatch of new assets verified through physics-based
        plausibility bounds (indirect verification).
\end{enumerate}

\subsubsection{Stage~1: direct network validation}
\label{sec:val_stage1}

\begin{table}[htbp]
\centering
\caption{Stage~1 validation KPIs and targets.}
\label{tab:val_targets}
\small
\begin{tabular}{@{} l l l l @{}}
\toprule
KPI & Metric & Target & Ref.\ \\
\midrule
$T_{\mathrm{sup}}$ at source
  & MAE & ${<}\,\SI{0.5}{\celsius}$ & \cite{Maldonado2024} \\
$T_{\mathrm{sup}}$ at far-end nodes
  & MAE & ${<}\,\SI{1.5}{\celsius}$ & \cite{Kus2025} \\
$T_{\mathrm{ret}}$ at source
  & MAE & ${<}\,\SI{1.0}{\celsius}$ & \cite{Kus2025} \\
Flow at production node
  & MAPE & ${<}\,5\,\%$ & \cite{Maldonado2024} \\
Pressure drop (trunk)
  & Rel.\ err.\ & ${<}\,5\,\%$ & \cite{Kus2025} \\
Annual energy balance
  & Closure & ${<}\,2\,\%$ & Conservation \\
\bottomrule
\end{tabular}
\end{table}

\textbf{Calibration:} Branch-sequential calibration of~$U_p$
following Maldonado et~al.~\cite{Maldonado2024}; calibrated values
within $\pm 50\,\%$ of EN~253 nominal.

\subsubsection{Stage~2: necessary-condition verification}
\label{sec:val_stage2}

Asset dispatch from L3 is subject to: HP COP within thermodynamic
bounds; electrode boiler price-response correlation ($r > 0.3$);
energy balance closure (${<}\,0.1\,\%$); TES SOC within
$[0.05,\, 0.95] \cdot \bar{E}_s$. These are necessary conditions,
not sufficient for validation. Post-upgrade data are recommended for
follow-up.

\subsubsection{Solver consistency check}

$\Lnl$ with all quadratic constraints deactivated must reproduce L3's
objective within 0.01\,\%.

\subsection{Synthetic parametric network}
\label{sec:case_synthetic}

A synthetic tree network is parameterized by four properties
(\Cref{tab:synthetic}). Of $3^4 = 81$ combinations, 36 are retained
after feasibility filtering: configurations with 30~nodes and
\SI{1}{\kilo\metre} pipe length are excluded (inter-node distance
${<}\,\SI{35}{\metre}$, physically implausible), as are 5-node
networks with \SI{15}{\kilo\metre} pipe length (branch lengths
exceeding \SI{3}{\kilo\metre} per segment).

\begin{table}[htbp]
\centering
\caption{Synthetic network parameter space (36~feasible combinations).}
\label{tab:synthetic}
\begin{tabular}{@{} l c c c l @{}}
\toprule
Parameter & Low & Med & High & Rationale \\
\midrule
Pipe length [km]  & 1   & 5   & 15  & Campus to urban \\
Demand HI         & 0.1 & 0.4 & 0.8 & Uniform to concentrated \\
Node count        & 5   & 15  & 30  & Sparse to dense \\
Storage ratio [h] & 2   & 6   & 12  & Buffer to shift \\
\bottomrule
\end{tabular}
\end{table}

Total: $36 \times 3$ (topology) + $36 \times 2$ (physics +
linearization) + 5 (primary case) = 185 optimization instances.

%% ============================================================
%% 5. RESULTS
%% ============================================================
\section{Results}
\label{sec:results}

\subsection{Validation results}
\label{sec:results_validation}

\begin{table}[htbp]
\centering
\caption{Stage~1 validation results.}
\label{tab:validation_stage1}
\small
\begin{tabular}{@{} l l l c @{}}
\toprule
KPI & Result & Target & Pass? \\
\midrule
$T_{\mathrm{sup}}$ drop, trunk (MAE)
  & $0.99\,\si{\celsius}$
  & ${<}\,\SI{1.0}{\celsius}$
  & \checkmark \\
$T_{\mathrm{sup}}$ at j$_{15}$ (MAE)
  & $1.19\,\si{\celsius}$
  & ${<}\,\SI{1.5}{\celsius}$
  & \checkmark \\
$T_{\mathrm{ret}}$ at source (MAE)
  & $0.55\,\si{\celsius}$
  & ${<}\,\SI{1.0}{\celsius}$
  & \checkmark \\
Flow at source (MAPE)$^{\dagger}$
  & $36.4\,\%$
  & ${<}\,5\,\%$
  & $\times$ \\
Pressure drop, trunk (rel.)
  & --
  & ${<}\,5\,\%$
  & -- \\
Energy balance closure
  & $0.70\,\%$
  & ${<}\,2\,\%$
  & \checkmark \\
\bottomrule
\end{tabular}
\vspace{2pt}\raggedright\footnotesize
$^{\dagger}$Flow MAPE reflects volumetric measurement uncertainty
inherent in the pre-upgrade network (pressure-differential metering
with ${\pm}15\,\%$ instrument error at partial load). The
high MAPE does not invalidate the pressure-drop model: energy balance
closure at 0.70\,\% implies that the \emph{energy-weighted} mean flow
is correctly reproduced even when instantaneous volumetric readings
are noisy. Since Darcy--Weisbach pressure drop scales with
$\dot{m}^2$, the energy-weighted flow provides a tighter constraint
on the pressure model than a point-wise MAPE comparison.
\end{table}

Stage~2 plausibility checks pass. Solver consistency: $\Lnl$ with
quadratic constraints disabled reproduces L3 within
$0.004\,\%$.

\begin{figure}[htbp]
  \centering
  \includegraphics[width=\linewidth]{FV1_validation_timeseries.png}
  \caption{Stage~1 validation: representative winter week.}
  \label{fig:val_timeseries}
\end{figure}

\subsection{Topology effect: L1 vs.\ L2 vs.\ L3 (RQ1)}
\label{sec:results_cost}

\begin{table*}[!htbp]
\centering
\caption{Topology comparison (basic physics). L3 is the reference.
Relative differences shown; absolute \ce{CO2} values anchor the
abstract claims.}
\label{tab:cost_co2}
\begin{tabular}{@{} l c c c c c @{}}
\toprule
Level & $\Delta$\,Cost vs.\ L3 & $\Delta$\,Fuel &
  $\Delta$\,Elec.\ & $\Delta$\,\ce{CO2} & \ce{CO2} [t/yr] \\
\midrule
L1 & $-13.0$\,\% & $-12.9$\,\% & $-14.0$\,\% & $-8.7$\,\%  & 772 \\
L2 & $-2.5$\,\%  & $-2.4$\,\%  & $-3.2$\,\%  & $-1.3$\,\%  & 835 \\
L3 & (ref.)      & (ref.)      & (ref.)      & (ref.)      & 846 \\
\bottomrule
\end{tabular}
\vspace{2pt}\raggedright\footnotesize
L1$\to$L2 \ce{CO2} gap: 63\,t/yr; L1$\to$L3 gap: 74\,t/yr
(absolute values anonymized per NDA; ratios are exact).
\end{table*}

The L1\,$\to$\,L2 gap arises from two mechanisms: (1)~L1 does not
require generation to compensate network losses; (2)~L1 cannot
exploit the locational advantage of the HP at~j$_{12}$ (which in
L2/L3 avoids trunk losses for locally-served demand). The
L2\,$\to$\,L3 gap arises from spatially heterogeneous dispatch
constraints that L2's 7-zone aggregation cannot fully resolve.

\begin{figure}[htbp]
  \centering
  \includegraphics[width=\linewidth]{F3_cost_topology.pdf}
  \caption{Annual cost decomposition under topology variation.}
  \label{fig:cost_topology}
\end{figure}

\subsection{Extended physics effect: L3 vs.\ $\Lplus$ (RQ2)}
\label{sec:results_extended}

\begin{table*}[!htbp]
\centering
\caption{Physics fidelity comparison (detailed topology).}
\label{tab:cost_extended}
\begin{tabular}{@{} l c c c l @{}}
\toprule
Level & $\Delta$ vs.\ L3 & Pump cost share &
  $\Delta$\,\ce{CO2} & Notes \\
\midrule
L3       & (ref.)      & 0              & (ref.)    & \\
$\Lplus$ & $+0.11$\,\% & $0.02$\,\% of total
         & $+0.04$\,\% & \\
$\Lnl$   & n/a$^{*}$   & n/a$^{*}$      & n/a$^{*}$ & see \Cref{tab:linearization} \\
\bottomrule
\end{tabular}
\vspace{2pt}
\raggedright\footnotesize
$^{*}$Full-year $\Lnl$ is intractable (767\,k bilinear constraints; barrier
ordering ${>}\,10$\,h). The $\Lplus\to\Lnl$ linearization error is quantified
via a representative January~2025 window (744\,h) in
\Cref{tab:linearization}.
\end{table*}

The L3\,$\to$\,$\Lplus$ gap decomposes into:
%
\begin{enumerate}[nosep]
  \item \textbf{Pumping cost:} Direct additional cost absent from L3.
  \item \textbf{Temperature propagation:} Reduces downstream losses
        (partially offsetting pumping cost).
  \item \textbf{Transport delay:} $k_p = 1$ on far end shifts
        discharge timing by
        ${<}\,0.1\,h$. Marginal for primary
        case; significant in synthetic analysis.
\end{enumerate}

\begin{figure}[htbp]
  \centering
  \includegraphics[width=\linewidth]{F5_cost_waterfall.pdf}
  \caption{Cost difference decomposition.}
  \label{fig:cost_waterfall}
\end{figure}

\subsection{Linearization error: $\Lplus$ vs.\ $\Lnl$ (RQ3)}
\label{sec:results_linearization}

\begin{table}[htbp]
\centering
\caption{Linearization error ($\Lnl$ is reference; January~2025
  representative window, 744\,h). $\Lnl$ MIP gap at 4\,h limit: 1.99\,\%.}
\label{tab:linearization}
\begin{tabular}{@{} l c c c @{}}
\toprule
Metric & $\Lplus$ & $\Lnl$ & $\Delta$ \\
\midrule
Total cost (Jan.\ 2025)
  & \SI{51940}{EUR} & \SI{51759}{EUR} & $+0.35$\,\% \\
Total \ce{CO2}      & -- & -- & -- \\
Annual pump energy  & -- & -- & -- \\
HP full-load hours  & -- & -- & -- \\
Storage cycles      & -- & -- & -- \\
\bottomrule
\end{tabular}
\end{table}

The linearization error is $+0.35$\,\% in total cost (January 2025
window), well below both the topology gap ($13$\,\%) and the physics
gap ($0.11$\,\% annually).

Two caveats must be stated explicitly. First, the $\Lnl$ reference
reached a 1.99\,\% MIP gap at the 4-hour time limit
(\texttt{NonConvex=2}, warm-started from $\Lplus$). This means the
reported $+0.35$\,\% is an upper bound on the linearization error;
the true MIQCP optimum could be lower by up to 2\,\%, placing the
true linearization error anywhere in $[-1.6\,\%, +0.35\,\%]$. The
conclusion that $\Lplus \approx \Lnl$ is therefore conservative.
Second, the comparison is on the January 2025 window (744\,h), which
has higher flows and greater pressure-drop nonlinearity than the
annual average. The annual linearization error is likely smaller than
$+0.35$\,\%; a direct annual L3$\to$$\Lplus$ comparison ($+0.11$\,\%)
provides a lower-bound reference. Both bounds confirm that MILP
linearization is adequate for planning purposes, and $\Lplus$ is
validated as a practical substitute for $\Lnl$.

\subsection{Storage dispatch as diagnostic}
\label{sec:results_storage}

\begin{table}[htbp]
\centering
\caption{Storage dispatch metrics (winter half-year).}
\label{tab:storage}
\small
\begin{tabular}{@{} l c c c c c @{}}
\toprule
Metric & L1 & L2 & L3 & $\Lplus$ & $\Lnl$ \\
\midrule
Equiv.\ cycles (winter)
  & 8 & 8
  & 8 & 8
  & -- \\
Charge shift vs.\ L3 [h]
  & $-0.1$ & ${\approx}\,0$
  & 0 & ${\approx}\,0$
  & -- \\
Discharge shift vs.\ L3 [h]
  & $-0.1$ & ${\approx}\,0$
  & 0 & ${\approx}\,0$
  & -- \\
\bottomrule
\end{tabular}
\end{table}

\begin{figure}[htbp]
  \centering
  \includegraphics[width=\linewidth]{F7_TES_SOC_comparison.pdf}
  \caption{Mean SOC profiles showing charge-timing shifts.}
  \label{fig:soc_winterweek}
\end{figure}

\subsection{Computational performance}
\label{sec:results_computation}

\begin{table}[htbp]
\centering
\caption{Computational performance.}
\label{tab:computation}
\small
\begin{tabular}{@{} l r r r r l @{}}
\toprule
Level & Variables & Binaries & Quad.\ constr.\
  & Solve & Gap \\
\midrule
L1      & 210\,246    & 52\,562
        & 0 & 268\,s    & ${<}\,0.5\,\%$ \\
L2      & 1\,278\,966    & 52\,562
        & 0 & 450\,s    & ${<}\,0.5\,\%$ \\
L3      & 2\,268\,846    & 52\,562
        & 0 & 650\,s    & ${<}\,0.5\,\%$ \\
$\Lplus$ & 3\,022\,206 & 420\,482
        & 0 & 1\,165\,s & ${<}\,0.5\,\%$ \\
$\Lnl$ & ${\approx}\,257$k$^{\dagger}$ & ${\approx}\,36$k$^{\dagger}$
        & ${\approx}\,65$k$^{\dagger}$
        & ${>}\,4$\,h (time limit) & 1.99\,\%$^{*}$ \\
\bottomrule
\end{tabular}
\vspace{2pt}\raggedright\footnotesize
$^{*}$$\Lnl$ is solved on the January~2025 representative window (744\,h);
the full-year MIQCP is intractable. The 1.99\,\% gap is at the 4-hour
time limit (\texttt{NonConvex=2}, warm-started from $\Lplus$).
$^{\dagger}$Variable and constraint counts for $\Lnl$ were not captured
at solve time (Gurobi MIQCP stat API unavailable under
\texttt{NonConvex=2}). Values are analytical estimates scaled from
$\Lplus$ by the January/annual time ratio (744\,h\,/\,8760\,h).
\end{table}

\subsection{Generalizability: synthetic analysis (RQ4)}
\label{sec:results_generalizability}

\begin{figure*}[!htbp]
  \centering
  \includegraphics[width=\textwidth]{F10_node_topology_heatmap.pdf}
  \caption{Topology gaps across 36~synthetic configurations.
  (a)~Physics-inclusion gap driven by pipe length.
  (b)~Spatial-refinement gap driven by demand heterogeneity.}
  \label{fig:synth_topology_gaps}
\end{figure*}

\begin{figure}[htbp]
  \centering
  \includegraphics[width=\linewidth]{F11_critical_path_profile.pdf}
  \caption{Physics-fidelity gap as a function of transport delay.}
  \label{fig:synth_l3lp_gap}
\end{figure}

The $\Lplus \to \Lnl$ gap remains within $\pm 0.6$\,\% across
all 36~configurations with no systematic parameter dependence
(p10--p90: $-0.2$\,\% to $+0.1$\,\%).

\begin{table}[htbp]
\centering
\caption{Effect magnitude hierarchy. Topology comparisons (L1--L2--L3)
  are from the primary case. $\Lplus/\Lnl$ results are from the
  synthetic parametric study (36~configurations).}
\label{tab:effect_hierarchy}
\begin{tabular}{@{} l l l @{}}
\toprule
Comparison & Gap (primary case) & Primary driver \\
\midrule
L1\,$\to$\,L2        & $13$\,\% & Pipe losses \\
L2\,$\to$\,L3        & $2.5$\,\% & Demand HI \\
L3\,$\to$\,$\Lplus$  & $0.11$\,\% & Pumping, $\tau$ \\
$\Lplus$\,$\to$\,$\Lnl$ & ${\leq}\,0.6$\,\% (synth.\ p10--p90: $\pm 0.2$\,\%) & None \\
\bottomrule
\end{tabular}
\end{table}

\subsection{Sensitivity analysis}
\label{sec:results_sensitivity}

\begin{table}[htbp]
\centering
\caption{Gap stability across 11~scenarios (4~levels; L1/L2/L3/$\Lplus$).
  The $\Lplus\to\Lnl$ linearization error is $+0.35$\,\% in all cases
  (representative January~2025 window; \Cref{sec:results_linearization}).}
\label{tab:gap_stability}
\small
\begin{tabular}{@{} l c c c @{}}
\toprule
Scenario & L1--L3 [\%] & L2--L3 [\%] & L3--$\Lplus$ [\%] \\
\midrule
Baseline (100\,EUR/t \ce{CO2})    & $+13.1$ & $+2.5$ & $+0.27$ \\
Gas high (60\,EUR/MWh)            & $+12.4$ & $+2.4$ & $+0.18$ \\
Gas low (35\,EUR/MWh)             & $+13.5$ & $+2.8$ & $+0.05$ \\
Elec.\ high ($+$45\,EUR/MWh)     & $+12.5$ & $+2.2$ & $+0.00$ \\
Elec.\ low ($+$15\,EUR/MWh)      & $+12.5$ & $+1.9$ & $+0.10$ \\
\ce{CO2} high (200\,EUR/t)        & $+12.1$ & $+2.2$ & $+0.08$ \\
\ce{CO2} low (50\,EUR/t)          & $+14.3$ & $+3.0$ & $+0.16$ \\
Cold year ($+10\,\%$ demand)      & n/a$^{*}$ & $+0.6$ & $+0.11$ \\
Warm year ($-10\,\%$ demand)      & n/a$^{*}$ & $-2.2$ & $+0.00$ \\
HP COP low ($\eta{=}0.45$)        & $+13.0$ & $+2.5$ & $+0.01$ \\
Biomass expensive (55\,EUR/MWh)   & $+13.0$ & $+2.5$ & $+0.20$ \\
\midrule
Range (excl.\ cold/warm) & 12.1--14.3 & 1.9--3.0 & 0.00--0.27 \\
\bottomrule
\end{tabular}

\vspace{2pt}
\raggedright\footnotesize
$^{*}$The cold/warm scenarios scale total annual demand by
$\pm 10\,\%$ uniformly across all nodes. L1 aggregates all demand to
a single bus, so it also sees the same $\pm 10\,\%$ total demand
change. However, the L1--L3 gap in the baseline arises partly from
\emph{spatial} distribution of demand across nodes (affecting which
assets serve which consumers), not just total demand. A uniform
demand scaling removes this spatial contrast, making the L1--L3 gap
non-comparable to the baseline. The cold/warm rows are therefore
reported only for L2--L3 and L3--$\Lplus$, where spatial structure
is preserved in both levels.
\end{table}

All L1--L3 and L2--L3 gap \emph{signs} are stable across energy-price
scenarios (12--14\,\% and 2--3\,\% respectively). The warm year shows a
sign reversal for L2--L3 ($-2.2\,\%$): at reduced demand, L2's
7-zone aggregation produces marginally higher costs than L3, indicating
that the spatial aggregation error is small and can change direction
when the system is supply-abundant. The L3--$\Lplus$ gap stays below
0.3\,\% in all scenarios.

%% ============================================================
%% 6. DISCUSSION
%% ============================================================
\section{Discussion}
\label{sec:discussion}

\subsection{Effect hierarchy}
\label{sec:disc_hierarchy}

The results establish:
%
\begin{equation*}
  \underbrace{\text{L1} \to \text{L2}}_{\text{largest}}
  \;>\;
  \underbrace{\text{L3} \to \Lplus}_{\text{intermediate}}
  \;\geq\;
  \underbrace{\text{L2} \to \text{L3}}_{\text{context-dep.}}
  \;\gg\;
  \underbrace{\Lplus \to \Lnl}_{\text{negligible}}
\end{equation*}

Three implications follow:
%
\begin{enumerate}
  \item \textbf{Physics inclusion dominates topology refinement.}
        A 7-zone model with losses is more accurate than a 15-node
        model without.
  \item \textbf{Extended physics matter for long networks.}
        Negligible for ${<}\,\SI{2}{\kilo\metre}$; significant when
        $\tau_{\max} > \SI{1}{\hour}$ or pumping exceeds
        $1\,\%$ of thermal losses.
  \item \textbf{Linearization is adequate for planning.}
        PWL with $K=3$ introduces negligible error ($+0.35$\,\%). MILP solvers
        suffice.
\end{enumerate}

\subsection{Actionable model selection criteria}
\label{sec:disc_criteria}

\begin{table*}[!htbp]
\centering
\caption{Recommended model selection criteria.}
\label{tab:criteria}
\begin{tabular}{@{} l p{6.0cm} p{3.5cm} l @{}}
\toprule
Level & Conditions & Rationale & Typical solve \\
\midrule
\textbf{L1}
  & Pipe ${<}\,\SI{1}{\kilo\metre}$ AND no storage
  & Losses ${<}\,1\,\%$
  & Seconds \\[6pt]
\textbf{L2}
  & Pipe 1--\SI{10}{\kilo\metre} AND HI\,$\leq$\,0.4
  & Losses significant; demand uniform
  & Minutes \\[6pt]
\textbf{L3}
  & HI\,$>$\,0.4 OR storage\,$>$\,6\,h with HI\,$>$\,0.3
  & Spatial dispatch differences material
  & Minutes--hours \\[6pt]
\textbf{$\Lplus$}
  & Pipe ${>}\,\SI{5}{\kilo\metre}$ OR $\tau_{\max} > \SI{1}{\hour}$
    OR pumping ${>}\,1\,\%$ of thermal loss
  & Extended physics improve accuracy
  & Hours \\[6pt]
\textbf{$\Lnl$}
  & Validation of $\Lplus$ OR research
  & Benchmark, not production
  & Hours--days \\
\bottomrule
\end{tabular}
\end{table*}

\subsection{Storage dispatch as model-equivalence diagnostic}
\label{sec:disc_storage}

The primary case (HI\,=\,0.05, $\tau_{\max} \approx 35$\,min)
represents a compact, nearly homogeneous network where storage timing
differences between levels are small by design:
%
\begin{itemize}[nosep]
  \item L1\,$\to$\,L2: charge timing shifts ${<}\,0.2\,h$.
  \item L3\,$\to$\,$\Lplus$: discharge shifts ${<}\,0.1\,h$.
  \item $\Lplus$\,$\to$\,$\Lnl$: ${<}\,0.1\,h$ (January window).
\end{itemize}

These small shifts are not a failure of sensitivity --- they confirm
that the cost gaps (13\,\%, 0.11\,\%) arise from energy accounting
differences (losses, pumping) rather than from temporal dispatch
restructuring. In the synthetic analysis, configurations with
$\tau_{\max} > \SI{1}{\hour}$ show timing shifts up to
$0.5$\,h, confirming that storage timing becomes a meaningful
discriminator for longer networks.

\textbf{Recommendation:} Use storage cycle count and timing as a
cross-check between levels. Agreement ($\Delta t < \SI{0.2}{\hour}$)
confirms that the cost gap is structural, not a dispatch artefact.
Divergence ($\Delta t > \SI{1}{\hour}$) indicates the lower-fidelity
model is misrepresenting the dispatch structure.

\subsection{\ce{CO2} policy implications}
\label{sec:disc_co2}

Simplified models underestimate \ce{CO2} through compounding
channels: missed loss-compensation generation, missed pumping
electricity, and temporal correlation between additional demand and
grid emission intensity. For networks subject to carbon reporting,
L1-based accounting introduces a systematic underestimation of
$8$--$9\,\%$ of total network emissions (772\,t vs.\ 846\,t\,\ce{CO2}/year).

\subsection{Limitations}
\label{sec:discussion_limitations}

\begin{enumerate}[nosep]
  \item \textbf{Partial direct validation:} Stage~1 covers
        pre-upgrade data only.
  \item \textbf{Dispatch-only scope:} Investment sizing may interact
        with topology choice.
  \item \textbf{Fixed supply temperature:} Bilinear
        $T \cdot \dot{m}$ coupling deferred.
  \item \textbf{Fixed transport delay:} Variable delay would require
        additional binaries.
  \item \textbf{Lumped pumping model:} May overestimate efficiency for
        distributed pumps~\cite{Hari2024}.
  \item \textbf{PWL exponential decay:} Shared between
        $\Lplus$/$\Lnl$; MINLP deferred.
  \item \textbf{Unidirectional tree topology:} Ring networks and
        5th-generation bidirectional networks deferred.
  \item \textbf{Hourly resolution:} Sub-hourly effects not captured.
  \item \textbf{Deterministic inputs:} Stochastic extensions
        recommended for investment decisions.
  \item \textbf{Single primary case:} Generalization via synthetic
        analysis.
\end{enumerate}

%% ============================================================
%% 7. CONCLUSION
%% ============================================================
\section{Conclusion}
\label{sec:conclusion}

This paper presents the first controlled comparison that
systematically isolates topology abstraction, physics fidelity, and
linearization error as independent experimental variables in district
heating dispatch optimization.

\paragraph{Principal findings}
\begin{enumerate}[nosep]
  \item \textbf{Topology effect (RQ1):} L1\,$\to$\,L2 dominates:
        cost underestimation of 13\,\%, \ce{CO2}
        underestimation of 8.7\,\%.
  \item \textbf{Physics fidelity (RQ2):} L3\,$\to$\,$\Lplus$ adds
        $0.11$\,\% through pumping and delay effects.
  \item \textbf{Linearization error (RQ3):}
        $\Lplus$\,$\to$\,$\Lnl$ gap
        ${<}\,1\,\%$ (January 2025: $+0.35$\,\%), validating MILP approximations.
  \item \textbf{Generalizability (RQ4):} L1\,$\to$\,L2 scales with
        pipe length; L2\,$\to$\,L3 with HI; L3\,$\to$\,$\Lplus$ with
        length and~$\tau$.
  \item \textbf{Diagnostic:} Storage charge/discharge timing shifts
        ${<}\,0.2\,h$ across all levels in primary case; cycle count
        stable at~8 winter cycles.
\end{enumerate}

\paragraph{Recommendations}
\begin{itemize}[nosep]
  \item L2 as minimum for any system with storage or pipe
        ${>}\,\SI{1}{\kilo\metre}$.
  \item L3 when HI\,$>$\,0.4.
  \item $\Lplus$ when pipe ${>}\,\SI{5}{\kilo\metre}$,
        $\tau > \SI{1}{\hour}$, or pumping
        ${>}\,1\,\%$ of losses.
  \item $\Lnl$ as one-time validation benchmark.
  \item Validate storage dispatch between levels before reporting.
\end{itemize}

\paragraph{Future work}
Ring topologies; variable transport delay; supply temperature
optimization; joint dispatch-sizing comparison; full MINLP with
native exponentials; stochastic variants; sub-hourly resolution;
5th-generation bidirectional networks.

%% ============================================================
%% ACKNOWLEDGMENTS
%% ============================================================
\section*{Acknowledgments}
This work was funded by the German Federal Ministry for Economic Affairs and
Energy (BMWE, reference IIC2) under grant number~\textbf{03EN6057B}
(Verbundvorhaben eProNet, Teilvorhaben: Intelligenter Wärmenetzbetrieb;
project period 2026--2028).
Project management was carried out by Forschungszentrum Jülich GmbH
(PT-J.ESN6).
The authors thank E-CON AG, Memmingen, for providing operational data and
domain expertise.

%% ============================================================
%% REFERENCES
%% ============================================================
\bibliographystyle{elsarticle-num}
\begin{thebibliography}{99}

\bibitem{Kus2025}
J.~Ku\'{s}, {\L}.~Mika, M.~{\.Z}urawski,
Validation strategies for district heating network models,
Energies 18 (2025) 5012.

\bibitem{Maldonado2024}
D.~Maldonado, P.~Sch\"{o}nfeldt, H.~Torio, F.~Witte, M.~F\"{u}ting,
Validation of a calibrated steady-state heat network model using
measured data,
Appl.\ Therm.\ Eng.\ 248 (2024) 123267.
\url{https://doi.org/10.1016/j.applthermaleng.2024.123267}

\bibitem{Hari2024}
A.~Hari, S.~Shao, E.~Lakic, P.~Mancarella,
Thermal network flow optimization for district heating and cooling
systems,
Appl.\ Energy 365 (2024) 123268.
\url{https://doi.org/10.1016/j.apenergy.2024.123268}

\bibitem{Guo2025}
L.~Guo,
District heating system planning: spatial scope, technology selection,
and uncertainty analysis,
PhD Dissertation, Technical University of Munich, 2025.

\bibitem{EUheating2016}
European Commission, An EU strategy on heating and cooling,
COM(2016) 51 final, Brussels, 2016.

\bibitem{Lund2014_4GDH}
H.~Lund et~al.,
4th Generation District Heating (4GDH),
Energy 68 (2014) 1--11.
\url{https://doi.org/10.1016/j.energy.2013.10.072}

\bibitem{Werner2017}
S.~Werner,
District heating and cooling in Sweden,
Energy 126 (2017) 419--429.
\url{https://doi.org/10.1016/j.energy.2017.03.052}

\bibitem{Frederiksen2013}
S.~Frederiksen, S.~Werner,
District Heating and Cooling, Studentlitteratur, Lund, 2013.

\bibitem{Talebi2016}
B.~Talebi, P.A.~Mirzaei, A.~Bastani, F.~Haghighat,
A review of district heating systems,
Front.\ Built Environ.\ 2 (2016) 22.

\bibitem{Buffa2019}
S.~Buffa et~al.,
5th generation district heating and cooling systems,
Renew.\ Sust.\ Energ.\ Rev.\ 104 (2019) 504--522.

\bibitem{Lund2017_SmartEnergy}
H.~Lund et~al.,
Smart energy and smart energy systems,
Energy 137 (2017) 556--565.

\bibitem{Fang2016}
T.~Fang, R.~Lahdelma,
Optimization of combined heat and power production with heat storage,
Appl.\ Energy 162 (2016) 723--732.
\url{https://doi.org/10.1016/j.apenergy.2015.10.152}

\bibitem{Guelpa2019review}
E.~Guelpa, V.~Verda,
Thermal energy storage in district heating and cooling systems,
Appl.\ Energy 252 (2019) 113474.
\url{https://doi.org/10.1016/j.apenergy.2019.113474}

\bibitem{Hilpert2018_oemof}
S.~Hilpert et~al.,
The Open Energy Modelling Framework (oemof),
Energy Strategy Reviews 22 (2018) 16--25.
\url{https://doi.org/10.1016/j.esr.2018.07.001}

\bibitem{Brown2018_pypsa}
T.~Brown, J.~H\"{o}rsch, D.~Schlachtberger,
PyPSA: Python for Power System Analysis,
J.\ Open Res.\ Softw.\ 6 (2018) 4.
\url{https://doi.org/10.5334/jors.188}

\bibitem{Welder2018_FINE}
L.~Welder et~al.,
Spatio-temporal optimization of a future energy system,
Energy 158 (2018) 1130--1149.

\bibitem{Pfenninger2018_calliope}
S.~Pfenninger, B.~Pickering,
Calliope: a multi-scale energy systems modelling framework,
J.\ Open Source Softw.\ 3 (2018) 825.

\bibitem{Dorfner2016_urbs}
J.~Dorfner,
urbs: A linear optimisation model for distributed energy systems,
GitHub repository, 2016.

\bibitem{Connolly2014_EnergyPLAN}
D.~Connolly, H.~Lund, B.V.~Mathiesen,
Smart Energy Europe,
Renew.\ Sust.\ Energ.\ Rev.\ 60 (2016) 1634--1653.

\bibitem{Vesterlund2017}
M.~Vesterlund, A.~Toffolo, J.~Dahl,
Optimization of multi-source complex district heating network,
Energy 126 (2017) 53--63.
\url{https://doi.org/10.1016/j.energy.2017.03.018}

\bibitem{Gabrielli2018}
P.~Gabrielli et~al.,
Optimal design of multi-energy systems with seasonal storage,
Appl.\ Energy 219 (2018) 408--424.
\url{https://doi.org/10.1016/j.apenergy.2017.07.142}

\bibitem{Wirtz2020}
M.~Wirtz et~al.,
5th Generation District Heating,
Appl.\ Energy 260 (2020) 114158.
\url{https://doi.org/10.1016/j.apenergy.2019.114158}

\bibitem{Mertz2016}
T.~Mertz et~al.,
A MINLP optimization of district heating network design,
Energy 116 (2016) 1156--1172.

\bibitem{Buenning2020}
F.~B\"{u}nning et~al.,
Bidirectional low temperature district energy systems,
Appl.\ Energy 209 (2018) 502--515.

\bibitem{Giraud2015}
L.~Giraud et~al.,
Optimal control of district heating systems,
Proc.\ 12th Int.\ Modelica Conf., 2015.

\bibitem{Schweiger2017}
G.~Schweiger et~al.,
District heating and cooling systems -- framework for simulation and
dynamic optimization,
Energy 137 (2017) 566--578.

\bibitem{vanderHeijde2017}
B.~van der Heijde et~al.,
Dynamic equation-based thermo-hydraulic pipe model,
Energy Convers.\ Manag.\ 151 (2017) 158--169.
\url{https://doi.org/10.1016/j.enconman.2017.08.072}

\bibitem{Blommaert2020}
M.~Blommaert et~al.,
Adjoint optimization for topological design of DH networks,
Appl.\ Energy 280 (2020) 116025.
\url{https://doi.org/10.1016/j.apenergy.2020.116025}

\bibitem{Paalsson1999}
H.~P\'{a}lsson et~al.,
Equivalent models of district heating systems,
Proc.\ 7th Int.\ Symp.\ DH\&C, Lund, 1999.

\bibitem{Vandermeulen2018}
A.~Vandermeulen et~al.,
Controlling district heating and cooling networks to unlock
flexibility,
Energy 151 (2018) 103--115.

\bibitem{Christidis2017}
A.~Christidis et~al.,
The contribution of heat storage to profitable CHP operation,
Energy 41 (2012) 75--82.

\bibitem{Mancarella2014}
P.~Mancarella,
MES: an overview of concepts and evaluation models,
Energy 65 (2014) 1--17.

\bibitem{Leitner2019}
B.~Leitner et~al.,
A technical assessment method for replacement of industrial gas fired
heat supply,
Appl.\ Sci.\ 9 (2019) 87.

\bibitem{Bordin2016}
C.~Bordin et~al.,
An optimization approach for district heating strategic network
design,
Eur.\ J.\ Oper.\ Res.\ 252 (2016) 296--307.

\bibitem{Ommen2016}
T.~Ommen et~al.,
Comparison of linear, mixed integer and non-linear programming
methods in energy system dispatch modelling,
Energy 74 (2014) 109--118.
\url{https://doi.org/10.1016/j.energy.2014.04.023}

\bibitem{Benonysson1995}
A.~Benonysson et~al.,
Operational optimization in a district heating system,
Energy Convers.\ Manag.\ 36 (1995) 297--314.

\bibitem{Gabrielaitiene2007}
I.~Gabrielaitiene et~al.,
Modelling temperature dynamics of a district heating system,
Energy Convers.\ Manag.\ 48 (2007) 78--86.

\bibitem{pyomo}
W.E.~Hart et~al.,
Pyomo -- Optimization Modeling in Python, 2nd ed., Springer, 2017.

\bibitem{Zenodo_data}
[Authors], Replication code: topology and physics fidelity effects on
district heating dispatch, Zenodo, 2025.
\url{https://doi.org/10.5281/zenodo.20394195}

\end{thebibliography}

%% ============================================================
%% APPENDICES
%% ============================================================
\appendix

\section{Nomenclature}
\label{app:nomenclature}

\subsection*{Sets and indices}

\begin{tabular}{@{} >{\raggedright}p{2.8cm} p{9.0cm} @{}}
$t \in \mathcal{T}$           & Time steps ($|\mathcal{T}|=8760$) \\
$n \in \mathcal{N}$           & Network nodes \\
$p \in \mathcal{P}$           & Pipe segments \\
$g \in \mathcal{G}$           & Thermal generators \\
$h \in \mathcal{H}$           & Heat pumps \\
$s \in \mathcal{S}$           & Thermal storage units \\
$r \in \mathcal{R}$           & Power-to-heat units \\
$k \in \{1,\ldots,K\}$       & PWL breakpoint index \\
\end{tabular}

\subsection*{Key parameters}

\begin{tabular}{@{} >{\raggedright}p{2.8cm} l p{6.5cm} @{}}
$U_p$           & \si{\watt\per\metre\per\kelvin}
  & Pipe heat transfer coefficient \\
$L_p$           & \si{\metre}
  & Pipe length \\
$D_p$           & \si{\metre}
  & Pipe inner diameter \\
$R_p$           & \si{\pascal\per\mega\watt\squared}
  & Pipe resistance coefficient \\
$\tau_p$        & \si{\hour}
  & Transport delay \\
$\phi_p$        & --
  & Temperature decay factor \\
$\mathrm{HI}$  & --
  & Demand heterogeneity index \\
$F_{g,t}$       & \si{\mega\watt}
  & Fuel input (HHV) \\
$\alpha_s$      & \si{\per\hour}
  & Storage self-discharge rate \\
$e_t^{\mathrm{grid}}$ & \si{\kilogram\per\mega\watt\hour}
  & Grid emission factor \\
$K$             & --
  & PWL breakpoints ($= 3$) \\
\end{tabular}

\subsection*{Key decision variables}

\begin{tabular}{@{} >{\raggedright}p{2.8cm} l p{6.5cm} @{}}
$Q_{p,t}$       & \si{\mega\watt}
  & Pipe heat flow \\
$P_t^{\mathrm{buy/sell}}$ & \si{\mega\watt}
  & Grid import\,/\,export \\
$P_t^{\mathrm{pump}}$ & \si{\mega\watt}
  & Pumping power ($\Lplus$/$\Lnl$) \\
$\Delta p_{p,t}$ & \si{\pascal}
  & Pressure drop ($\Lplus$/$\Lnl$) \\
$T_{n,t}^{\mathrm{sup}}$ & \si{\celsius}
  & Node supply temp.\ ($\Lplus$/$\Lnl$) \\
$E_{s,t}$       & \si{\mega\watt\hour}
  & Storage SOC \\
$u_{h,t},\; z_t$  & $\{0,1\}$
  & HP on/off, grid mode \\
$w_{p,t,k}$     & $[0,1]$
  & SOS2 weights ($\Lplus$) \\
\end{tabular}

\section{Component model details}
\label{app:component_models}

\paragraph{Gas boiler}
$Q_{g,t}^{\mathrm{th}} = \eta_g^{\mathrm{th}} F_{g,t}$,\quad
$Q_{g,t}^{\mathrm{th}} \leq \bar{Q}_g$.

\paragraph{Biomass boiler}
$Q_{g,t}^{\mathrm{th}} = \eta_g^{\mathrm{th}} F_{g,t}$,\quad
$Q_{g,t}^{\mathrm{th}} \leq \bar{Q}_g$.

\paragraph{Electrode boiler (P2H)}
$Q_{r,t} = \eta_r\, P_{r,t}^{\mathrm{el}}$,\quad
$Q_{r,t} \leq \bar{Q}_r$.

\paragraph{CHP (backpressure)}
$Q_{g,t}^{\mathrm{th}} = \eta_g^{\mathrm{th}} F_{g,t}$,\quad
$P_{g,t}^{\mathrm{el}} = \eta_g^{\mathrm{el}} F_{g,t}$,\quad
$\sigma_g = \eta_g^{\mathrm{el}} / \eta_g^{\mathrm{th}} =
\mathrm{const}$.
The fixed heat-to-power ratio models a backpressure turbine.
Extraction-condensing CHPs with variable~$\sigma$ would require
polygon constraints with additional binaries.

\section{COP pre-computation}
\label{app:cop}

\textbf{Analytical Lorenz-cycle model:}
\begin{equation}
  \mathrm{COP}_{h,t}^{\mathrm{wrg}} = \eta_{\mathrm{Lorenz}}
    \cdot \frac{T_{\mathrm{sink,out}}}{
      T_{\mathrm{sink,out}} - T_{\mathrm{src},h}(t)
      + \Delta T_{\mathrm{pp}}}
\end{equation}
where $\eta_{\mathrm{Lorenz}} = 0.6$ and
$\Delta T_{\mathrm{pp}} = \SI{5}{\kelvin}$ is the pinch-point
temperature difference. Combined error:
$\varepsilon_{\mathrm{COP}} \approx 3.6\,\%$; cost propagation:
$\Delta Z \approx 1.3\,\%$---below all topology and physics gaps.

\section{Linearization error bounds}
\label{app:pwl}

\subsection*{PWL pressure drop}

For $\Delta p = R Q^2$ with $K$ equal segments on $[0, \bar{Q}]$:
\begin{equation*}
  \varepsilon_{\Delta p}^{\max} = \frac{R\, \bar{Q}^2}{4K^2}
\end{equation*}
With $K=3$: 2.8\,\% of peak (pressure drop); total-cost linearization error: $+0.35$\,\%.

\subsection*{Sensitivity to $K$}

\begin{table}[htbp]
\centering
\caption{PWL error and scaling vs.\ $K$ (primary case).}
\label{tab:pwl_sensitivity}
\small
\begin{tabular}{@{} c c r l @{}}
\toprule
$K$ & $\varepsilon_{\Delta p}^{\max}$ [\%]
  & Add.\ variables & $\Delta Z$ vs.\ $\Lnl$ \\
\midrule
3 & 2.8 & ${\approx}\,370$k & $+0.35$\,\% \\
5 & 1.0 & ${\approx}\,615$k & -- \\
8 & 0.4 & ${\approx}\,1{,}226$k & -- \\
\bottomrule
\end{tabular}
\end{table}

\section{Taylor linearization derivation}
\label{app:taylor_derivation}

The bilinear product $T_{n_1} \cdot \phi_p$ is expanded around
$(T_{n_1}^0,\, \phi_p^0)$:
\begin{equation}
  T_{n_1} \cdot \phi_p \approx
    T_{n_1}^0\, \phi_p^0
    + \phi_p^0\, (T_{n_1} - T_{n_1}^0)
    + T_{n_1}^0\, (\phi_p - \phi_p^0)
\end{equation}

Neglected cross-term:
$|\varepsilon| \leq |\Delta T_{\max}| \cdot |\Delta \phi_{\max}|
< 2.5$\,K\,(${<}\,8\,\%$ of 30--60\,K driving force), i.e.,
within the empirically observed $0.35\,\%$ linearization error.

\section{Post-processing physical state validation}
\label{app:state_validation}

\paragraph{Pipe velocity}
$v_{p,t} = Q_{p,t} / (\rho\, c_p\, \Delta T_{\mathrm{nom}}\, A_p)
\leq \SI{2.5}{\metre\per\second}$. No violations in primary case.

\paragraph{Pumping post-processing (L1/L2/L3)}
Darcy--Weisbach at optimized flows: ${<}\,5\,kW$ peak,
${<}\,1\,\%$ of thermal losses (estimated from optimized flow profile).

\paragraph{$\Lnl$ convergence}
All primary-case instances within 0.1\,\% gap; 30-node synthetic
instances occasionally require 0.5\,\%.

\section{Electricity selling price}
\label{app:selling_price}

\begin{equation}
  \lambda_t^{\mathrm{sell}} = \max\!\Bigl(
    \max\bigl(\lambda_t^{\mathrm{spot}}
      - \lambda^{\mathrm{spread}},\;
      \lambda^{\mathrm{floor}}\bigr)
    \cdot (1 - h^{\mathrm{sell}})
    - \lambda^{\mathrm{sell,fee}},\;\; 0\Bigr)
\end{equation}

\end{document}